\documentclass[stu, floatsintext, 12pt, spanish]{apa7}

% ==========================================
% Paquetes de Idioma, Codificación y Fuentes
% ==========================================
\usepackage[utf8]{inputenc}
\usepackage[spanish, es-tabla]{babel}
\usepackage{csquotes}

% ==========================================
% Paquetes Gráficos y Ruta de Imágenes
% ==========================================
\usepackage{graphicx}
\graphicspath{{figuras/}} % Búsqueda automática en la carpeta figuras/
\usepackage{hyperref}
\usepackage{amssymb} 

% ==========================================
% Paquete para cadenas de búsqueda
% ==========================================
\usepackage{listings}
\lstdefinestyle{cadenas}{
	basicstyle=\ttfamily\small,
	breaklines=true,
	columns=fullflexible,
	keepspaces=true,
	frame=none,
	xleftmargin=1.27cm,
	aboveskip=0.5em,
	belowskip=0.5em,
}

% ==========================================
% Configuración Bibliográfica (APA 7 - BibLaTeX + Biber)
% ==========================================
\usepackage[backend=biber, style=apa]{biblatex}
\addbibresource{referencias.bib}
% BibLaTeX permite cortes en números y letras de identificadores URL extensos.
\setcounter{biburlnumpenalty}{100}
\setcounter{biburllcpenalty}{100}
\setcounter{biburlucpenalty}{100}
% Mantener cada referencia completa en una página; permitir saltos entre entradas.
\appto\bibsetup{\interlinepenalty=10000\relax}

% ==========================================
% Justificar texto
% ==========================================
\usepackage{ragged2e}
\justifying

% ==========================================
% Configuración de sangría
% ==========================================
\usepackage{indentfirst}
\setlength{\parindent}{1.27cm}


%Estilos tablas

\usepackage{tabularx}   
\usepackage{longtable}
\usepackage{booktabs}   
\usepackage{array}     
\usepackage{ragged2e}  
\usepackage{float}

% ==========================================
% Configuración para diagrama de Gantt
% ==========================================
\usepackage{pgfgantt}
\usepackage{xcolor}


\definecolor{completado}{RGB}{112,173,71}
\definecolor{proceso}{RGB}{244,177,131}
\definecolor{grislinea}{RGB}{217,217,217}

% ==========================================
% Estilos Personalizados UNHSJM
% ==========================================
\usepackage{unhsjm-formato}

% Ruta del Logo Institucional (ubicado en figuras/)
\unhsjmlogo{figuras/logo-unhsjm.jpeg}

% ==========================================
% Portada e Información Institucional (UNHSJM)
% ==========================================
\universidad{UNIVERSIDAD NACIONAL HÉROES SAN JOSÉ DE LAS MULAS}
\facultad{FACULTAD DE CIENCIAS E INGENIERÍAS}
\carrera{CARRERA DE INGENIERÍA EN SISTEMAS DE INFORMACIÓN}

% Modalidad de Culminación de Estudios (Cambiar según corresponda: "PROYECTO DE GRADO" o "MONOGRAFÍA")
\modalidad{PROYECTO DE GRADO}
% \modalidad{MONOGRAFÍA}

% Título y Subtítulo del Trabajo
\title{Travel DLS}
\subtitulo{Plataforma Web para la gestión, localización y seguimiento de carga pesada}
\shorttitle{Sistema Web TravelDLS}

% Autores (Trabajo Grupal o Individual)
% Para agregar múltiples integrantes, separe las líneas con \\
\autores{
  Estudiante 1 - Jareth Said Acevedo Bonilla (Carnet: 2240351) \\
  Estudiante 2 - Ashly Deney Castro Velazquez (Carnet: 2240097) \\
  Estudiante 3 - Elsa Lineth Cerda Ugarte (Carnet: 2240337)
}

% Tutor / Asesor del Trabajo
\tutor{M.Sc. Wilfredo José Meneses Bejarano}

% Ubicación y Fecha
\ciudad{Managua, Nicaragua}
\fecha{Septiembre, 2026}

\begin{document}

% Evitar saltos de página inmediatamente después de una palabra cortada.
\brokenpenalty=10000\relax

\renewcommand{\listtablename}{Índice de tablas}
\renewcommand{\listfigurename}{Índice de figuras}

% Generar Portada Institucional
\maketitle

% Tabla de Contenidos / Índice
\tableofcontents
\clearpage
\phantomsection
\addcontentsline{toc}{section}{Índice de tablas}
\listoftables
\clearpage
\phantomsection
\addcontentsline{toc}{section}{Índice de figuras}
\listoffigures
\clearpage

% ==========================================
% Inclusión de Capítulos
% ==========================================
% Capítulo 1: Introducción
% Propósito: Describir los antecedentes, planteamiento del problema, justificación, formulación y los objetivos (general y específicos) del proyecto de grado.


\section{Introducción}

El transporte terrestre de carga forma parte del movimiento diario de mercancías entre empresas, ciudades y diferentes puntos de distribución. Con el paso del tiempo, la forma de administrar estas operaciones también ha cambiado, especialmente por la incorporación de tecnologías que permiten manejar una mayor cantidad de información durante los procesos logísticos. En este sentido, la logística inteligente ha comenzado a apoyarse en herramientas digitales para automatizar actividades, monitorear envíos y mejorar la organización de las operaciones \parencite{kalkha2023rising}.

Este avance tecnológico ha hecho que la gestión del transporte ya no se observe únicamente desde el traslado físico de una mercancía. Actualmente también se estudian aspectos relacionados con el seguimiento de los envíos y la planificación de los recorridos. \textcite{zhang2024ecommerce}, por ejemplo, desarrollaron una propuesta en la que el seguimiento de paquetes y la planificación de rutas forman parte de una misma solución logística. Por otro lado, \textcite{xin2022logistics} muestran que la búsqueda de recorridos más adecuados continúa siendo un tema de interés dentro de la distribución de mercancías, debido a su relación con el uso de los vehículos, los tiempos de recorrido y la organización de las entregas.

Tomando como referencia esta evolución de los sistemas logísticos surge TravelDLS, un proyecto orientado al transporte terrestre de carga pesada. La propuesta consiste en desarrollar una plataforma web que permita reunir diferentes empresas transportistas dentro de un mismo entorno y facilitar la relación entre empresas, conductores y clientes. Además de servir como espacio para administrar las operaciones propias de cada empresa, se busca que las empresas transportistas puedan dar a conocer sus servicios y tener mayor visibilidad ante personas o negocios que necesiten contratar transporte de carga. El sistema también incorporará funciones relacionadas con los viajes, la ubicación de las unidades y la planificación de recorridos, manteniendo cada una de estas funciones conectada con el servicio que se está realizando.

A partir de esta propuesta, el presente proyecto se encuentra organizado en diferentes apartados que permiten desarrollar de manera progresiva la propuesta de TravelDLS. En primer lugar, se presenta el \textbf{\hyperref[sec:planteamiento]{Planteamiento del Proyecto}}, donde se describe la situación que da origen a la propuesta y las necesidades que se pretenden atender. Posteriormente, se establecen los \textbf{\hyperref[sec:objetivos]{Objetivos}} que orientan el desarrollo del proyecto y se presentan los \textbf{\hyperref[sec:antecedentes]{Antecedentes}}, los cuales permiten conocer investigaciones y soluciones relacionadas con la gestión del transporte, el seguimiento y la optimización de rutas. A continuación, la \textbf{\hyperref[sec:justificacion]{Justificación}} expone la importancia de la propuesta, mientras que el \textbf{\hyperref[sec:alcance]{Alcance}} y las \textbf{\hyperref[sec:limitaciones]{Limitaciones}} delimitan las funciones que serán desarrolladas y los aspectos que quedan fuera de esta etapa. Seguidamente, se analiza la \textbf{\hyperref[sec:viabilidad]{Viabilidad del Proyecto}} desde las perspectivas técnica, operativa, económica y legal. El \textbf{\hyperref[sec:marco-teorico]{Marco Teórico}} reúne los antecedentes del proyecto, los conceptos y las bases teóricas necesarias para comprender los componentes de TravelDLS. Finalmente, la \textbf{\hyperref[sec:metodologia]{Metodología}} describe el procedimiento seguido para la selección de investigaciones, la recolección y análisis de información y la organización del proceso de desarrollo de la plataforma.

\subsection{Planteamiento del Proyecto}
\label{sec:planteamiento}

El transporte terrestre de carga pesada implica mucho más que mover una mercancía de un punto a otro. En cada servicio también deben coordinarse vehículos, conductores, cargas, recorridos, tiempos y comunicación con el cliente. Cuando esa información se encuentra separada o no se actualiza de forma oportuna, resulta más difícil saber qué está ocurriendo durante el viaje y reaccionar ante cambios, atrasos o inconvenientes. Esta necesidad de mantener una mejor comunicación también aparece en estudios sobre logística, donde se señalan dificultades entre clientes, conductores y empresas, además de la necesidad de conocer el estado de los envíos durante el traslado \parencite[p. 2]{sultan2022real}.

El seguimiento y la planificación de rutas son dos apoyos que pueden ayudar a ordenar este tipo de operaciones. \textcite{pauzuoliene2024smart} encontraron que, en una muestra de 46 empresas logísticas de Lituania, el 69.6 \% utilizaba tecnologías para dar seguimiento a la carga en tiempo real y el 60.9 \% empleaba programas para planificar o seleccionar rutas (p. 988). En Centroamérica también se han señalado necesidades de digitalización y coordinación. En Panamá, por ejemplo, se han identificado barreras tecnológicas, baja digitalización y dificultades de coordinación dentro del transporte terrestre \parencite[p. 10]{castillo2025potenciando}. Aunque estos estudios corresponden a otros contextos, ayudan a entender que mantener conectada la información del viaje sigue siendo una necesidad presente en distintas operaciones de transporte.

En Nicaragua ya existen trabajos que atienden partes de esta situación. \textcite{aragon2018implementacion} desarrollaron un sistema para administrar fletes terrestres con la intención de ordenar los controles y facilitar el acceso a la información logística (p. 2). Por su parte, \textcite{garcia2022desarrollo} abordaron en León la necesidad de conocer recorridos y ubicación de unidades de transporte mediante una aplicación con funciones de geolocalización en tiempo real (pp. 3, 7). Son proyectos diferentes a TravelDLS, pero muestran que tanto la organización de la información como la ubicación de los vehículos ya han sido necesidades abordadas desde la informática en el país.

A partir de este panorama, TravelDLS se concentra en una necesidad más específica: que las empresas dedicadas al transporte terrestre de carga pesada puedan organizar sus operaciones y, al mismo tiempo, mantener conectada la información que necesitan los conductores y los clientes durante el servicio. El problema no consiste únicamente en registrar un viaje. También implica relacionar el vehículo, el conductor, la carga, la ruta, la ubicación y el estado del servicio para que la información tenga sentido durante todo el recorrido.

Por ello, la pregunta que guía el proyecto es la siguiente: ¿Cómo mejorar la gestión, coordinación, localización, seguimiento y planificación de rutas de las operaciones de transporte terrestre de carga pesada, de manera que empresas transportistas, conductores y clientes dispongan de información centralizada y oportuna durante el servicio?

Como respuesta a esta necesidad se propone TravelDLS, una plataforma web bajo el modelo de Software como Servicio (SaaS). La idea es reunir en un mismo entorno la información de empresas transportistas, vehículos, conductores, cargas, viajes, rutas y clientes, de manera que las distintas partes del servicio no queden aisladas. La plataforma también contempla la publicación y contratación de servicios, el seguimiento de los viajes, la consulta del estado de la carga y un módulo para apoyar la planificación de rutas con varios destinos. TravelDLS no pretende reemplazar las decisiones de la empresa, sino ofrecer un medio más ordenado para consultar información y apoyar la gestión diaria de los viajes.

\subsection{Objetivos}
\label{sec:objetivos}

\subsubsection{Objetivo General}
 Desarrollar un sistema web denominado TravelDLS para la gestión, localización y trazabilidad en tiempo real de las operaciones de transporte terrestre de carga pesada, incorporando la optimización matemática de rutas para fortalecer los procesos administrativos y logísticos de las empresas transportistas y brindar mayor transparencia a los clientes durante el traslado de sus cargas.

\subsubsection{Objetivos Específicos}
\label{sec:objetivos-especificos}
\begin{itemize}
	\item Analizar los requerimientos funcionales y no funcionales del sistema, evaluando y comparando métodos de optimización de rutas basados en el Problema del Agente Viajero (Vecino Más Cercano, 2-opt y Google OR-Tools) para seleccionar la solución más eficiente.
	\item Diseñar la arquitectura del sistema, incluyendo la base de datos y la estructura de los módulos que conformarán la plataforma.
	\item Proponer un prototipo de plataforma web (backend y frontend), implementando el módulo de optimización con Google OR-Tools para la gestión eficiente de usuarios, vehículos, conductores, cargas, viajes y rutas.
\end{itemize}

\subsection{Antecedentes}
\label{sec:antecedentes}

Para entender mejor de dónde parte TravelDLS, se revisaron investigaciones realizadas de manera Internacional como Nacional que abordan problemas parecidos desde distintos puntos de vista. Algunos trabajos se concentran en el seguimiento de la carga, otros en la planificación de rutas y otros en la administración de los servicios de transporte. Revisarlos permite reconocer qué soluciones se han probado, qué resultados han obtenido y qué ideas pueden servir como referencia para el proyecto, sin asumir que una propuesta aplicada en otro contexto funcionará de la misma manera en TravelDLS.

\phantomsection\label{ant:pauzuoliene}
En el ámbito internacional, \textcite{pauzuoliene2024smart} realizaron la investigación titulada \textbf{Smart technologies integration and challenges in the context of logistics companies}, en donde estudiaron el uso de tecnologías inteligentes en empresas logísticas de Lituania. La investigación contó con 46 representantes del sector y mostró que el 69.6 \% utilizaba seguimiento de carga en tiempo real y el 60.9 \% empleaba programas para planificar o seleccionar rutas. Los autores también relacionan estas herramientas con la posibilidad de conocer la ubicación, el estado de la carga y los tiempos esperados de entrega (pp. 988-989). Para TravelDLS, este trabajo es útil porque muestra que el seguimiento y la planificación pueden formar parte de una misma gestión logística, aunque el estudio se enfoca en la adopción de tecnologías y no en el desarrollo de una plataforma para carga pesada.

\phantomsection\label{ant:sultan}
También a nivel internacional, \textcite{sultan2022real} llevaron a cabo el estudio titulado \textbf{A real-time mobile and web-based logistics management system for collaborative urban freight transportation} en el cual propusieron un sistema móvil y web para apoyar el transporte urbano de mercancías. La solución reúne funciones de registro, reservas, conductores, administración, inventario, seguimiento de pedidos y planificación de rutas (pp. 1-3). Los autores además señalan problemas de comunicación entre clientes, conductores y empresas, junto con la necesidad de que el cliente pueda conocer el estado de su envío durante el traslado (p. 2). Este trabajo se acerca a TravelDLS por la forma en que integra varios participantes y procesos dentro de una misma solución, aunque TravelDLS estará dirigido específicamente al transporte terrestre de carga pesada y a varias empresas usuarias.

\phantomsection\label{ant:aguilar}
Al acercar la revisión a Centroamérica, \textcite{aguilar2022optimizacion}, en Honduras,publicaron el artículo titulado \textbf{Optimización de Rutas de compañías distribuidoras y surtidoras en Tegucigalpa} en donde se estudió la optimización de rutas para compañías distribuidoras y surtidoras de Tegucigalpa. En el escenario analizado, el recorrido total pasó de 7,037.8 km a 4,008 km después de utilizar software de optimización y se estimó una reducción del 30 \% en los costos de distribución (p. 56). Más que tomar ese porcentaje como una meta para TravelDLS, el estudio sirve para mostrar que una propuesta de rutas puede evaluarse con datos concretos, como la distancia recorrida y los costos asociados.

\phantomsection\label{ant:hernandez}
Otro trabajo realizado en Honduras por \textcite{hernandez2024evaluacion} ejecutaron el proyecto de investigación titulado \textbf{Evaluación de alternativas de software de planificación de rutas para la empresa transporte Gonzales} en el cual se comparó distintas alternativas de software para la planificación de rutas en la empresa Transporte Gonzales. El estudio evaluó herramientas como Route4Me, Geotab y DispatchTrack y reportó una reducción calculada del 17 \% en los costos totales, menor que la meta inicial del 20 \% (pp. 10, 111). A la vez, la opción seleccionada permitía visualizar rutas, estado de envíos y tiempos estimados de entrega (p. 111). Este resultado es importante para TravelDLS porque recuerda que no conviene asumir de antemano que una herramienta será la mejor: la comparación debe hacerse con criterios claros y tomando en cuenta el contexto en el que se va a utilizar.

En Panamá, \textcite{castillo2025potenciando} se encontró una investigación titulada \textbf{Potenciando la competitividad logística de Panamá: una ingeniería inversa del transporte terrestre frente a Singapur y Rotterdam} analizó el transporte terrestre desde una perspectiva de competitividad logística e identificó, entre otros aspectos, costos elevados, limitaciones de infraestructura, baja digitalización y dificultades de coordinación (p. 10). Aunque el trabajo no desarrolla una aplicación informática, ayuda a entender que los problemas del transporte no dependen únicamente de la tecnología. TravelDLS puede apoyar la organización, el seguimiento y la planificación de los viajes, pero no pretende resolver por sí solo todos los factores que afectan al sector.

\phantomsection\label{ant:aragon}
En Nicaragua, \textcite{aragon2018implementacion} presentaron la investigación titulada \textbf{Implementación de un sistema de información para la administración de fletes terrestres en el departamento de transporte de la empresa master logistics, s.a}, cuyo propósito fue administrar fletes terrestres en dicha empresa. El trabajo buscó ordenar controles, mejorar la exactitud de los datos y facilitar el acceso a la información de las actividades logísticas; además, su objetivo general fue implementar un sistema capaz de registrar los elementos que forman parte de la administración de fletes (pp. 2, 9). Este antecedente es cercano a TravelDLS por su relación con el transporte de carga, aunque la propuesta actual amplía ese enfoque al incorporar contratación, seguimiento y planificación de rutas dentro de una plataforma pensada para distintas empresas.

\phantomsection\label{ant:garcia}
También en Nicaragua, \textcite{garcia2022desarrollo} presentaron la investigación titulada \textbf{Desarrollo de una aplicación móvil Android para el seguimiento y localización de las unidades de transporte público en la zona urbana del municipio de León, marzo - diciembre, año 2021} en la UNAN-León, la cual consistió en una aplicación Android para el seguimiento y localización de unidades de transporte público. El problema estudiado estuvo relacionado con la falta de información sobre recorridos, paradas y ubicación, y uno de sus objetivos fue mostrar en tiempo real la localización de las unidades mediante servicios de Google Maps (pp. 3, 7). Aunque se trata de transporte de pasajeros y no de carga pesada, el trabajo aporta una referencia nacional para el componente de ubicación de TravelDLS y demuestra que este tipo de tecnología ya se ha aplicado en proyectos de transporte desarrollados en el país.

Al observar los trabajos en conjunto, se aprecia que la administración del transporte, el seguimiento y la planificación de rutas han sido tratados desde diferentes enfoques. Los estudios internacionales muestran experiencias donde seguimiento y planificación forman parte de soluciones logísticas, mientras que los trabajos regionales permiten comparar resultados de optimización que no siempre son iguales. En Nicaragua se encontraron propuestas para administrar fletes y para localizar unidades de transporte, pero esas capacidades aparecen en proyectos separados. Con la revisión realizada hasta este momento, TravelDLS encuentra una oportunidad en reunir varias de estas funciones dentro de una plataforma orientada a la carga pesada. Esta idea todavía debe comprobarse con el mapeo sistemático completo, por lo que se mantiene como una brecha preliminar y no como una conclusión definitiva.

\subsection{Justificación}
\label{sec:justificacion}

TravelDLS se justifica porque las actividades de una operación de carga están relacionadas entre sí. Registrar un viaje tiene poco valor si luego esa información no puede conectarse con el vehículo, el conductor, la carga, la ruta y el estado del servicio. La propuesta busca que esos datos puedan consultarse de una manera más ordenada durante el recorrido y que cada actor tenga acceso a la información que le corresponde. La importancia de estas funciones también se observa en empresas logísticas que ya utilizan herramientas de seguimiento y planificación de rutas como parte de sus operaciones \parencite[p. 988]{pauzuoliene2024smart}.

El módulo de optimización se incorpora como un apoyo para planificar los recorridos, no como una promesa de que un algoritmo siempre encontrará la mejor opción para cualquier situación. Los antecedentes regionales muestran precisamente que los resultados cambian según los datos y la herramienta empleada: \textcite{aguilar2022optimizacion} obtuvo una reducción estimada del 30 \% en el escenario estudiado, mientras que \textcite{hernandez2024evaluacion} reportaron un 17 \% en su evaluación (p. 56; p. 111). Por eso, en TravelDLS será importante comparar los métodos definidos en el proyecto antes de valorar su comportamiento.

Sobre esta base, la innovación propuesta consiste en adaptar la relación entre contratación, gestión y seguimiento al transporte de carga pesada en Nicaragua. Frente a los antecedentes nacionales de administración de fletes para una empresa y localización de transporte de pasajeros \parencite{aragon2018implementacion,garcia2022desarrollo}, TravelDLS vinculará el servicio contratado con los recursos asignados por cada transportista y la consulta del estado de la carga por el cliente. La mejora esperada es dar continuidad a la información desde la contratación hasta el seguimiento del viaje. Como ya existen sistemas logísticos integrados y plataformas multitenant \parencite{sultan2022real,lasluisa2026diseno}, el aporte no se atribuye a esas tecnologías, sino a su aplicación al proceso descrito, cuya utilidad se evaluará en el proyecto. Los \textbf{\hyperref[sec:antecedentes-proyecto]{Antecedentes del Proyecto}} precisan las diferencias con cada estudio.

Para sostener esta propuesta, la plataforma contempla un modelo de pago en el que el porcentaje destinado al mantenimiento se agregará cuando la empresa transportista establezca el precio del viaje, de forma que el cliente vea desde el inicio el monto final antes de aceptar el servicio. Lo recaudado servirá para apoyar el funcionamiento de la plataforma, cubrir correcciones, mejoras y otros servicios necesarios para mantenerla disponible.

A partir de estas razones, resulta necesario delimitar de manera concreta qué componentes y funcionalidades formarán parte de TravelDLS durante esta etapa del proyecto. Por ello, a continuación se establece el alcance de la propuesta, definiendo hasta dónde llegará el desarrollo del sistema y qué elementos serán considerados dentro de su implementación.

\subsection{Alcance}
\label{sec:alcance}

TravelDLS se desarrollará como una plataforma web bajo el modelo SaaS, pensada para que distintas empresas de transporte de carga pesada puedan utilizarla. En esta primera etapa se contempla:

\begin{itemize}
	\item Registro y gestión de perfiles para empresas transportistas, conductores y clientes.
	\item Administración de vehículos, conductores, cargas, viajes, rutas y servicios asociados a cada empresa.
	\item Publicación de servicios de transporte y consulta o contratación por parte de los clientes.
	\item Asignación de conductores y vehículos a los viajes, con consulta del estado de cada operación.
	\item Localización y seguimiento de los viajes mediante servicios de mapas y funciones de ubicación disponibles en los dispositivos utilizados.
	\item Consulta para el cliente del estado y ubicación de su carga durante el desarrollo del servicio.
	\item Comparación de los resultados de Vecino Más Cercano y 2-opt con Google OR-Tools, siendo esta última la herramienta que se implementará en TravelDLS.
	\item Integración de una pasarela de pago en la que el porcentaje destinado al mantenimiento forme parte del precio final mostrado al cliente antes de aceptar el servicio.
	\item Realización de pruebas funcionales y elaboración de la documentación técnica y los manuales de usuario.
\end{itemize}

Aunque estas funcionalidades permiten establecer lo que se pretende desarrollar en TravelDLS, el proyecto también debe reconocer aquellos aspectos que no serán cubiertos durante esta primera etapa. Definir estos límites permite mantener el trabajo dentro de las condiciones académicas, técnicas y temporales disponibles.

\subsection{Limitaciones}
\label{sec:limitaciones}

Como TravelDLS se desarrollará en el marco de este proyecto académico y se concentrará en sus funciones principales, será necesario establecer algunos límites para mantener el trabajo enfocado y realizable:

\begin{itemize}
	\item La solución se desarrollará como aplicación web; en esta etapa no se contempla una aplicación móvil nativa independiente.
	\item El proyecto estará enfocado en el transporte terrestre de carga pesada y no abarcará transporte aéreo, marítimo, ferroviario ni transporte público de pasajeros.
	\item La optimización se probará en viajes con varios destinos y con los métodos definidos para el proyecto, por lo que no se pretende cubrir todos los tipos de problemas de rutas que existen.
	\item La actualización de la ubicación dependerá de la conexión a internet, del dispositivo utilizado y de los servicios de ubicación disponibles.
	\item Las pruebas se realizarán con los escenarios disponibles dentro del proyecto académico, por lo que sus resultados no representan todavía una evaluación de uso a gran escala.
	\item En esta etapa no se incluirán conexiones especializadas con sistemas externos como plataformas contables, aduaneras, aseguradoras u otros servicios que requieran desarrollos adicionales.
	\item TravelDLS se plantea para ser utilizado por distintas empresas transportistas y no como una solución exclusiva para una sola organización.
\end{itemize}

\subsection{Viabilidad del proyecto}
\label{sec:viabilidad}

Una vez definido lo que formará parte del proyecto y los límites bajo los cuales será desarrollado, corresponde analizar si la propuesta puede llevarse a la práctica con los recursos, tecnologías y condiciones disponibles. Para ello, la viabilidad de TravelDLS se estudia desde diferentes perspectivas.

\subsubsection{Viabilidad Técnica}

TravelDLS presenta viabilidad técnica debido a que las tecnologías seleccionadas permiten desarrollar una plataforma web con una separación clara entre la interfaz de usuario, la lógica del servidor, el almacenamiento de la información, la gestión de archivos multimedia y la visualización geográfica generada por empresas transportistas, conductores y clientes. Para el frontend se utilizará Next.js, mientras que el backend será desarrollado con AdonisJS. La información será almacenada en PostgreSQL, las fotografías de usuarios y empresas y el logotipo de la aplicación se gestionarán mediante Cloudinary, y los mapas, rutas y seguimiento de ubicación se apoyarán en Mapbox. Además, se utilizarán JSON Web Tokens (JWT) para apoyar los procesos de autenticación y control de acceso. El frontend será desplegado en Vercel y el backend en Render, utilizando un dominio adquirido mediante GoDaddy o Namecheap y conexiones HTTPS protegidas mediante certificados SSL/TLS.

Como referencia para la organización interna del backend, se tomó en cuenta el trabajo de \textcite{arce2026arquitectura}, quienes desarrollaron una plataforma web utilizando una estructura basada en Modelo-Vista-Controlador (MVC) complementada con patrones como Service Layer, con el propósito de separar responsabilidades entre los distintos componentes del sistema y facilitar su mantenimiento (pp.~28, 33--35). Este antecedente resulta útil para TravelDLS porque muestra una forma de mantener la lógica de negocio separada de los controladores. Sin embargo, la arquitectura de TravelDLS no busca reproducir completamente la estructura utilizada en dicho trabajo, sino aplicar específicamente MVC + Service Layer dentro del backend desarrollado con AdonisJS.

En el backend de TravelDLS se implementará una estructura basada en el patrón Modelo-Vista-Controlador (MVC), complementada con una Service Layer. Esta organización se aplicará internamente en AdonisJS y no implica que Next.js forme parte de la Vista del MVC del backend. El Modelo representará las entidades y relaciones persistidas en PostgreSQL, mientras que los Controladores recibirán las solicitudes provenientes del frontend y coordinarán las respuestas de la API. La lógica principal del negocio se trasladará a la capa de servicios, donde se concentrarán procesos como la gestión de viajes, la asignación de vehículos y conductores, la contratación de servicios, el seguimiento y la optimización de rutas. En una plataforma web estudiada por \textcite{arce2026arquitectura}, la combinación de MVC con patrones como Service Layer permitió separar responsabilidades y favorecer la mantenibilidad y escalabilidad del sistema. A partir de este enfoque, TravelDLS utilizará MVC + Service Layer en el backend para mantener los controladores ligeros y separar las reglas del negocio de la gestión de solicitudes.

Next.js se utilizará como aplicación frontend independiente encargada de la interfaz de usuario y de consumir la API expuesta por AdonisJS. Esta separación permitirá desarrollar vistas específicas para empresas transportistas, conductores y clientes sin mezclar la lógica de presentación con las reglas del negocio. AdonisJS, por su parte, concentrará la lógica del servidor y facilitará la organización de módulos relacionados con usuarios, empresas, vehículos, cargas, viajes, rutas, servicios y pagos.

PostgreSQL se selecciona porque la información de TravelDLS posee una naturaleza principalmente relacional. Las empresas estarán asociadas con usuarios, vehículos, conductores, cargas y viajes, mientras que cada viaje se relacionará con rutas, estados y clientes. Por ello, una base de datos relacional resulta adecuada para conservar estas dependencias y mantener consistencia entre los registros. JWT se utilizará como mecanismo de apoyo para la autenticación entre el frontend y el backend, complementado con controles de autorización que permitan diferenciar los permisos de cada tipo de usuario.

Las fotografías de los usuarios y de las empresas, así como el logotipo de TravelDLS, se alojarán en Cloudinary para facilitar su carga, organización y visualización dentro de la plataforma. En PostgreSQL se conservarán las referencias necesarias para asociar cada imagen con su usuario, empresa o elemento de la aplicación. Esta integración tiene antecedentes en TaxiGo, que utiliza Cloudinary para gestionar archivos multimedia \parencite{yaqoub2024taxigo}, y en Truxi, donde las imágenes se almacenan en ese servicio y sus enlaces se relacionan con los registros del sistema \parencite{vad2020truxi}. Para los mapas se utilizará Mapbox, que permitirá representar las rutas y la ubicación de los vehículos durante los viajes. Su aplicación en soluciones de transporte y logística se observa en el sistema de seguimiento en tiempo real desarrollado por \textcite{aronsson2021tracking}, donde Mapbox se seleccionó como proveedor de mapas, y en Thunder Logistics, que integra Mapbox GL JS, Directions API y Geocoding API para mapas interactivos, rutas y seguimiento público \parencite{deyheralde2026thunder}. La posición se obtendrá de los servicios de ubicación del dispositivo del conductor, se enviará al backend y se actualizará en el mapa conforme se reciban nuevos datos; de esta manera, el seguimiento podrá mostrar la ubicación más reciente disponible.

Para el despliegue, Vercel alojará el frontend desarrollado en Next.js y Render alojará el backend construido con AdonisJS. Esta distribución permitirá actualizar ambos componentes de forma independiente y mantener la aplicación disponible desde una infraestructura centralizada. El dominio será adquirido mediante GoDaddy o Namecheap y se utilizará HTTPS mediante certificados SSL/TLS para proteger la comunicación entre los usuarios y los servicios del sistema.

En conjunto, la propuesta se considera técnicamente viable porque las tecnologías seleccionadas permiten separar la interfaz, la lógica del servidor y la persistencia de datos de acuerdo con las necesidades definidas para TravelDLS. La adopción de MVC + Service Layer específicamente en el backend permitirá ordenar el código y separar responsabilidades, mientras que la división entre frontend, backend y base de datos facilitará el mantenimiento y futuras ampliaciones de la plataforma.

Contar con tecnologías adecuadas y una estructura definida permite establecer que el sistema puede construirse desde el punto de vista técnico. Sin embargo, también es necesario analizar cómo estas herramientas podrán utilizarse dentro de las actividades diarias de empresas, conductores y clientes, lo que conduce al estudio de la viabilidad operativa.

\subsubsection{Viabilidad Operativa}

La viabilidad operativa de TravelDLS se basa en que será una plataforma web accesible mediante un navegador, por lo que las empresas transportistas no necesitarán instalar una versión independiente del sistema ni mantener servidores propios. El acceso estará dividido principalmente entre empresas transportistas, conductores y clientes, mostrando a cada usuario las funciones relacionadas con su participación dentro del servicio.

La utilización de sistemas digitales con diferentes tipos de usuarios cuenta con antecedentes dentro del área logística. \textcite{sultan2022real} desarrollaron una solución móvil y web que integra clientes, conductores y administradores mediante funciones de registro, reservas, seguimiento y gestión logística (pp.~1--4). Este antecedente es relevante para TravelDLS porque demuestra que los procesos de transporte pueden organizarse mediante una plataforma donde varios actores interactúan con información relacionada con un mismo servicio.

La utilidad operativa del seguimiento y la planificación también ha sido documentada en empresas logísticas. \textcite{pauzuoliene2024smart} encontraron que, dentro de la muestra analizada, el 69.6\,\% de las empresas utilizaba tecnologías para seguimiento de carga en tiempo real y el 60.9\,\% empleaba herramientas para planificación o selección de rutas (pp.~988--989). Aunque estos resultados corresponden a un contexto distinto al nicaragüense, muestran que ambas funciones forman parte de procesos tecnológicos utilizados en el sector logístico.

En TravelDLS, las empresas podrán administrar vehículos, conductores, cargas, viajes y servicios; los conductores podrán consultar la información de los recorridos asignados; y los clientes podrán conocer el estado de los servicios contratados. Al mantenerse la aplicación centralizada, las actualizaciones realizadas sobre la plataforma estarán disponibles para todos los usuarios sin requerir instalaciones adicionales en cada empresa.

Como parte de esa gestión, los usuarios y las empresas podrán incorporar sus fotografías a los perfiles y registros que les correspondan, mientras que el logotipo de TravelDLS se mostrará de forma consistente en la interfaz. La experiencia de TaxiGo muestra la utilidad de integrar la gestión de imágenes y las funciones de seguimiento dentro de un sistema de transporte con distintos tipos de usuarios \parencite{yaqoub2024taxigo}. Durante un viaje activo, el conductor compartirá su ubicación mediante el dispositivo que utilice y Mapbox la presentará sobre el mapa para que la empresa y el cliente consulten el avance del servicio según sus permisos. Truxi ofrece un antecedente de uso de mapas para indicar lugares de recogida y entrega y apoyar los recorridos de los conductores \parencite{vad2020truxi}; además, Chalmers documenta el uso de Mapbox dentro de un sistema de seguimiento posicional en tiempo real para una empresa de transporte \parencite{aronsson2021tracking}, y Thunder Logistics lo emplea para visualizar rutas, envíos activos y páginas de tracking con ubicación actualizada \parencite{deyheralde2026thunder}. En TravelDLS, esta información se vinculará con el viaje registrado, de modo que el seguimiento y la consulta de la ruta formen parte del trabajo cotidiano de los participantes.

El funcionamiento de TravelDLS dependerá de factores como la conexión a Internet, la disponibilidad de los servicios de ubicación y los dispositivos utilizados por los conductores. Por ello, estos elementos se considerarán durante el diseño de la plataforma, tomando como referencia investigaciones previas y documentación técnica relacionada con la digitalización de los procesos logísticos. Asimismo, la pertinencia operativa de la solución se analizará mediante los procedimientos de diseño y evaluación establecidos en la metodología DSRM, con el propósito de determinar si las funcionalidades propuestas responden al problema identificado y a los requerimientos definidos para el sistema.

Además de comprobar que TravelDLS puede integrarse a las actividades de sus posibles usuarios, es necesario determinar si su desarrollo y funcionamiento pueden sostenerse económicamente. Por esta razón, el siguiente análisis considera la inversión inicial, los costos operativos y los ingresos que podrían generarse mediante el modelo propuesto para la plataforma.

\subsubsection{Viabilidad Económica}

La inversión inicial de TravelDLS se calcula sumando el desarrollo del sistema, el dominio del primer año y la configuración, las pruebas y la puesta en marcha: C\$52,732.80 + C\$439.44 + C\$3,662.00 = C\$56,834.24. Para expresar en córdobas los importes presupuestados en dólares, se multiplica cada monto por C\$36.62. Este valor corresponde al redondeo a dos decimales del tipo de cambio oficial de C\$36.6243 por USD~1.00 publicado por el Banco Central de Nicaragua para 2026 \parencite{bcn2026tipocambio}. A continuación se explica la procedencia de los rubros que integran esta estimación académica.

Para el desarrollo del sistema se asigna una valoración académica de C\$17,577.60 por integrante. Al multiplicar este monto por los tres integrantes se obtiene C\$52,732.80. Como referencias salariales, World Salaries reporta un promedio mensual de C\$18,158 para desarrolladores y programadores en Nicaragua, mientras que WageIndicator presenta un rango mensual inicial de C\$15,612 a C\$31,324 \parencite{worldsalaries2026developer,wageindicator2026desarrolladores}. Estas referencias orientan la valoración del trabajo técnico; el monto adoptado por integrante es una estimación académica por dedicación parcial al proyecto. Se contabiliza como parte de la inversión por el valor del trabajo aportado, aunque no se desembolse como salario.

Para el dominio del primer año se presupuestan USD~12.00, equivalentes a C\$439.44 al multiplicar $12.00 \times 36.62$. Las referencias de Namecheap y GoDaddy permiten consultar precios de registro y renovación, que varían según el dominio, las promociones y el período contratado \parencite{namecheapdomain,godaddy2026domain}. Por ello, los USD~12.00 se utilizan como una asignación presupuestaria aproximada y no como una cotización definitiva. La configuración, las pruebas y la puesta en marcha se valoran internamente en USD~100.00, equivalentes a C\$3,662.00 mediante $100.00 \times 36.62$; este rubro corresponde a una estimación del equipo para preparar el funcionamiento de la plataforma.

GitHub Free y los certificados SSL/TLS de Let's Encrypt se incluyen con un costo inicial de C\$0.00, de acuerdo con sus servicios gratuitos publicados \parencite{githubpricing,letsencrypt2021about}. Asimismo, el presupuesto inicial contempla Cloudinary para las imágenes y Mapbox para los mapas dentro de sus niveles gratuitos \parencite{cloudinarypricing,mapboxpricing}. El costo de C\$0.00 se mantiene mientras el uso permanezca dentro de los límites aplicables de cada servicio. La Tabla~\ref{tab:inversion-inicial} reúne los rubros de la inversión inicial.

\begin{table}[H]
\centering
\caption{Inversión inicial estimada de TravelDLS}
\label{tab:inversion-inicial}
\begin{tabularx}{\textwidth}{>{\RaggedRight\arraybackslash}X r}
\toprule
Concepto & Monto (C\$) \\
\midrule
Desarrollo del sistema (tres integrantes) & 52,732.80 \\
Dominio del primer año & 439.44 \\
Configuración, pruebas y puesta en marcha & 3,662.00 \\
GitHub Free y certificado SSL/TLS de Let's Encrypt & 0.00 \\
Cloudinary y Mapbox dentro de sus niveles gratuitos & 0.00 \\
\midrule
Total & 56,834.24 \\
\bottomrule
\end{tabularx}
\end{table}

Los costos operativos se calculan a partir del alojamiento de la aplicación, la base de datos, el soporte y una reserva para consumos adicionales. Para el backend se toma como referencia una instancia de Render de USD~7.00 mensuales, equivalentes a C\$256.34 mediante $7.00 \times 36.62$ \parencite{renderpricing}. Para PostgreSQL se conserva una asignación presupuestaria de USD~10.00 mensuales, equivalente a C\$366.20; su costo definitivo deberá confirmarse según la capacidad y el almacenamiento contratados en Render. Para el frontend se considera la tarifa base de Vercel Pro de USD~20.00 mensuales, equivalente a C\$732.40 mediante $20.00 \times 36.62$, bajo el supuesto de una sola plaza de desarrollador de pago \parencite{vercelpricing}. El soporte y mantenimiento se estiman internamente en USD~100.00 mensuales, equivalentes a C\$3,662.00, para atender ajustes y tareas de seguimiento técnico.

Aunque se prevé iniciar con los niveles gratuitos de Mapbox y Cloudinary, el uso de mapas, solicitudes, almacenamiento y procesamiento de imágenes puede generar cargos al superar los límites de cada servicio \parencite{mapboxpricing,cloudinarypricing}. Para contemplar esa posibilidad, la Tabla~\ref{tab:costos-operativos} incluye una \textit{Reserva para consumos adicionales} de C\$732.40 mensuales, equivalente a USD~20.00. Su valor anual se obtiene multiplicando C\$732.40 por 12 meses, con un resultado de C\$8,788.80. Esta reserva ya forma parte del total de costos operativos y está destinada a atender posibles cargos adicionales o contribuir a la contratación de un plan de pago. Su suficiencia dependerá del consumo real y de las tarifas aplicables; si el costo supera lo reservado, deberá ajustarse el presupuesto.

La suma de los rubros mensuales es C\$256.34 + C\$366.20 + C\$732.40 + C\$3,662.00 + C\$732.40 = C\$5,749.34. Al multiplicar este total por 12 se obtienen C\$68,992.08 anuales. La Tabla~\ref{tab:costos-operativos} presenta el desglose mensual y anual, incluida la reserva.

\begin{table}[H]
\centering
\caption{Costos operativos estimados de TravelDLS}
\label{tab:costos-operativos}
\begin{tabularx}{\textwidth}{>{\RaggedRight\arraybackslash}X r r}
\toprule
Concepto & Mensual (C\$) & Anual (C\$) \\
\midrule
Backend en Render & 256.34 & 3,076.08 \\
PostgreSQL en Render & 366.20 & 4,394.40 \\
Frontend en Vercel Pro & 732.40 & 8,788.80 \\
Soporte y mantenimiento & 3,662.00 & 43,944.00 \\
Reserva para consumos adicionales & 732.40 & 8,788.80 \\
\midrule
Total & 5,749.34 & 68,992.08 \\
\bottomrule
\end{tabularx}
\end{table}

El dominio del primer año se incluye en la inversión inicial para no duplicarlo como gasto operativo. A partir del segundo año, el modelo agrega una renovación anual estimada de C\$439.44, por lo que los costos anuales proyectados ascienden a C\$69,431.52. Esta renovación conserva el supuesto de USD~12.00 del presupuesto, aunque las tarifas de renovación pueden diferir de las de registro y deberán confirmarse con el proveedor \parencite{namecheapdomain,godaddy2026domain}. La Tabla~\ref{tab:supuestos-economicos} resume los valores utilizados como base del modelo.

\begin{table}[H]
\centering
\caption{Supuestos base del modelo económico}
\label{tab:supuestos-economicos}
\begin{tabularx}{\textwidth}{>{\RaggedRight\arraybackslash}X r}
\toprule
Concepto & Valor (C\$) \\
\midrule
Inversión inicial & 56,834.24 \\
Costos operativos anuales & 68,992.08 \\
Renovación de dominio anual & 439.44 \\
\bottomrule
\end{tabularx}
\end{table}

Con estos costos como base, la viabilidad económica se proyecta a cinco años. El modelo de ingresos considera un recargo del 8\,\% sobre el precio definido por la empresa transportista. Para un envío base de C\$3,662.00, el cliente pagaría C\$3,954.96; de ese monto, C\$292.96 corresponden al ingreso bruto de TravelDLS. Al descontar una comisión estimada del 3\,\% por pasarela de pago, equivalente a C\$118.65, queda un ingreso neto de C\$174.31 por envío. Estos porcentajes son supuestos del modelo académico. El detalle del cálculo por envío se presenta en la Tabla~\ref{tab:anexo-contribucion-envio} del Anexo~\ref{anexo:flujo-caja}.

La proyección considera 60 envíos mensuales en el primer año y un aumento de 30 envíos mensuales por cada año posterior, hasta llegar a 180 envíos mensuales en el quinto año. El flujo completo se presenta en la Tabla~\ref{tab:anexo-flujo-caja-proyectado} del Anexo~\ref{anexo:flujo-caja}.

Con una tasa académica de descuento del 12\,\%, el modelo obtiene un valor actual neto (VAN) de C\$547,106.97, una tasa interna de retorno (TIR) de 162\,\% y un período de recuperación de la inversión (PRI) de 1.002 años. Estos indicadores se muestran en la Tabla~\ref{tab:anexo-indicadores-financieros} del Anexo~\ref{anexo:flujo-caja}. Bajo estos supuestos, TravelDLS presenta viabilidad económica preliminar, aunque los resultados deberán revisarse cuando se conozcan la comisión definitiva de pago, la demanda real, la renovación del dominio y el consumo efectivo de infraestructura.

No obstante, la sostenibilidad económica no constituye el único aspecto que debe considerarse antes de implementar TravelDLS. Debido a que la plataforma gestionará datos personales, ubicaciones, contrataciones y servicios relacionados con el transporte terrestre, también debe revisarse el marco legal aplicable a su funcionamiento.

\subsubsection{Viabilidad Legal}

La viabilidad legal de TravelDLS se sustenta en las normas nicaragüenses que regulan la protección de datos, el comercio electrónico y el transporte terrestre de carga. La Ley No. 787, Ley de Protección de Datos Personales, es necesaria porque la plataforma manejará información de empresas, conductores y clientes, incluyendo ubicaciones en tiempo real, fotografías de perfil y documentos de identidad. Su artículo 4 condiciona la creación de ficheros de datos personales a que exista autorización y consentimiento del titular, lo cual obliga a TravelDLS a implementar mecanismos de consentimiento informado durante el registro de usuarios. El artículo 5 pide que los datos sean los necesarios para el fin que se persigue, de modo que la plataforma no puede recolectar información más allá de la que realmente utiliza para gestionar los viajes. El artículo 7 obliga a informar al titular para qué se usarán sus datos, quién los verá y cómo puede ejercer sus derechos de acceso, rectificación o cancelación \parencite[Arts.4,5 y 7]{ley787_2012}. Sin esta ley, TravelDLS no tendría un marco claro para tratar la información de los usuarios, y las empresas transportistas podrían exponerse a reclamos por el mal uso de datos de sus conductores o clientes.

La contratación de servicios a través de TravelDLS también tiene respaldo legal. La Ley No. 729, Ley de Firma Electrónica, es necesaria porque la plataforma está diseñada para que empresas transportistas publiquen servicios y los clientes los contraten directamente desde la web, sin necesidad de firmar documentos en físico. Su artículo 1 otorga eficacia y valor jurídico a la firma electrónica y a los certificados digitales, así como a toda información en formato electrónico. El artículo 6 establece que la firma electrónica certificada tendrá el mismo valor jurídico que la firma manuscrita y será admisible como medio de prueba en procesos judiciales o administrativos \parencite[Arts. 1 y 6]{ley729_2010}. Esto significa que los contratos celebrados a través de la plataforma entre empresas, conductores y clientes tendrán pleno valor jurídico, y permiten que la contratación dentro de TravelDLS no dependa de trámites presenciales. La Ley 729 fue publicada en La Gaceta No. 165 del 30 de agosto de 2010, y es la norma vigente que regula esta materia en Nicaragua.

En cuanto a la seguridad de la información, el Código Penal, Ley No. 641, es necesario porque TravelDLS operará como una plataforma en línea que almacena información crítica para la operación de las empresas transportistas. Su artículo 246 castiga con prisión de uno a tres años a quien, con intención de dañar, distribuya programas informáticos destructivos, y su artículo 245 sanciona la destrucción o inutilización de registros informáticos, con penas mayores si se trata de información ligada a un servicio público \parencite[Arts.245--246]{ley641_2008}. Esto obliga a TravelDLS a proteger su infraestructura con medidas como los tokens JWT y las conexiones HTTPS que ya se mencionan en la propuesta técnica, y justifica la inversión en seguridad para evitar que un ataque informático deje a las empresas sin acceso a sus viajes, rutas o datos de clientes.

En el ámbito del transporte, la Ley No. 524, Ley General de Transporte Terrestre, y su Reglamento (Decreto No. 42-2005), son necesarios porque TravelDLS está dirigido específicamente al transporte terrestre de carga pesada. El artículo 5 de la Ley 524 clasifica el transporte de carga como el traslado de bienes y todo tipo de objetos, sólidos o líquidos, que se presta a otros. El artículo 7 exige que todos los vehículos de servicio público de transporte terrestre cumplan las normas técnicas de inspección y obtengan la constancia de aprobación correspondiente \parencite[Arts. 5 y 7]{ley524_2005}. El Reglamento, en su artículo 13, detalla que la inspección técnica mecánica es el proceso mediante el cual una unidad de transporte público se somete a revisión de funcionamiento y estado de sus partes, y debe realizarse por lo menos cada seis meses \parencite[Art.13]{decreto42_2005}. Aunque TravelDLS no es una empresa transportista, sus módulos de registro de vehículos, conductores y viajes deben alinearse con estos requisitos para que la plataforma realmente sirva de apoyo a las empresas en el cumplimiento de la normativa, y no se convierta en un sistema que ignore las obligaciones legales del sector.

La Ley No. 842, Ley de Protección de los Derechos de las Personas Consumidoras y Usuarias, también tiene relevancia. Es necesaria porque los clientes que contraten transporte de carga a través de TravelDLS podrían quedar en una situación de desventaja informativa frente a la empresa transportista. Su artículo 2 cubre todos los actos de consumo dentro del país e incluso protege a microempresas cuando hay asimetría informativa con el proveedor. El artículo 3 excluye a quienes compran para su actividad empresarial, pero los clientes que contraten servicios de carga podrían no tener el mismo conocimiento técnico que la empresa sobre rutas, tiempos o costos \parencite[Arts. 2--3]{ley842_2013}. Por eso, la plataforma debe mostrar de forma clara el precio final, incluyendo el recargo del 8 \% y las condiciones del servicio, para que el cliente sepa exactamente qué está contratando antes de aceptar.

Estas leyes no solo dan respaldo jurídico a TravelDLS, sino que definen cómo debe operar la plataforma en aspectos concretos como el consentimiento de datos, la validez de los contratos electrónicos, la seguridad de la información, el cumplimiento normativo del transporte y la protección del cliente. Sin este marco, la plataforma podría funcionar técnicamente, pero quedaría expuesta a reclamos legales y no sería una herramienta confiable para las empresas transportistas que la utilicen.

Una vez consideradas las condiciones técnicas, operativas, económicas y legales que rodean la propuesta, también resulta necesario observar las condiciones en las que los posibles usuarios interactuarán con la plataforma. El funcionamiento de TravelDLS no dependerá únicamente del cumplimiento de requisitos técnicos o normativos, sino también de factores relacionados con la adopción de tecnología, las formas habituales de trabajo, la comunicación entre los participantes, las condiciones de acceso digital y la confianza en la información proporcionada. Por esta razón, a continuación se analizan algunos aspectos socioculturales y tecnológicos que deben tomarse en cuenta durante el diseño de la plataforma.

\subsection{Aspectos socioculturales}
\label{sec:aspectos-socioculturales}

El desarrollo de TravelDLS no depende únicamente de la implementación de sus componentes tecnológicos, sino también de las condiciones bajo las cuales las empresas transportistas, conductores y clientes podrían utilizar la plataforma. En las MIPYMES latinoamericanas, la transformación digital presenta diferentes niveles de avance según el tamaño de la empresa, los recursos disponibles y su grado de madurez tecnológica \parencite{dini2021transformacion}. Por esta razón, aspectos como la adopción tecnológica, los hábitos de trabajo, la comunicación, el acceso digital y la confianza deben considerarse durante el diseño de una solución dirigida al transporte terrestre de carga pesada.

\subsubsection{Adopción de la tecnología}
\label{sec:adopcion-tecnologia}

La incorporación de nuevas herramientas digitales puede representar un cambio importante para las MIPYMES, especialmente cuando modifica procedimientos que los trabajadores ya realizan de una manera conocida. \textcite{dini2021transformacion} señalan que la digitalización de las pequeñas empresas no ocurre de forma homogénea, debido a diferencias en sus capacidades, recursos y nivel de desarrollo tecnológico. En Nicaragua, \textcite{barreda2023economia} encontraron que MIPYMES de Estelí ya utilizan dispositivos y herramientas digitales como computadoras, teléfonos inteligentes y redes sociales dentro de sus actividades, lo que muestra una incorporación progresiva de la tecnología en pequeñas empresas del país.

Ante este contexto, TravelDLS busca facilitar su adopción mediante una plataforma web con procesos organizados y una interfaz comprensible para sus distintos usuarios. Las funciones de registro de vehículos, conductores, cargas, viajes y rutas deben mantener una estructura consistente, evitando que la digitalización introduzca una complejidad mayor que los procedimientos que pretende mejorar. De esta manera, la usabilidad y una curva de aprendizaje razonable se consideran elementos importantes para favorecer la integración progresiva de la plataforma en las actividades de las empresas transportistas.

\subsubsection{Hábitos de trabajo}
\label{sec:habitos-trabajo}

La digitalización de los procesos logísticos no consiste únicamente en sustituir registros físicos por formularios electrónicos, sino en mejorar la forma en que la información se organiza y se relaciona durante una operación. \textcite{aleman2021sistema}, en un estudio desarrollado para la compañía Logística de Carga Intermodal en Managua, identificaron que la empresa utilizaba diferentes herramientas computacionales para procesos como cotizaciones, inventario y facturación, pero estas no satisfacían completamente sus necesidades, lo que motivó el desarrollo de un sistema integrado de gestión.

Esta situación guarda relación con TravelDLS, cuyo propósito es evitar que la información de una operación se gestione como elementos independientes. La plataforma busca relacionar la empresa, el cliente, la carga, el vehículo, el conductor, el viaje y la ruta dentro de un mismo servicio. Ante esta necesidad, TravelDLS responde mediante la centralización de información y la reutilización de datos previamente registrados, permitiendo que las actividades digitales sigan el flujo de trabajo de la operación y reduzcan la necesidad de mantener registros separados.

\subsubsection{Comunicación entre los actores}
\label{sec:comunicacion-actores}

Las operaciones de transporte requieren coordinación entre participantes que normalmente se encuentran en lugares diferentes. La empresa administra el servicio, el conductor realiza el recorrido y el cliente necesita conocer el avance de su carga. \textcite{heinbach2022datadriven} muestran que las plataformas digitales aplicadas al transporte terrestre pueden integrar información relacionada con órdenes, recursos y seguimiento, facilitando la visibilidad de las operaciones y el intercambio de información entre los participantes del proceso.

Ante esta realidad, TravelDLS busca disminuir la dependencia de comunicaciones individuales para consultar información básica del servicio. La empresa podrá conocer el estado de sus operaciones y recursos, el conductor tendrá acceso a la información correspondiente al viaje asignado y el cliente podrá consultar el estado y ubicación de su carga. Esto no pretende reemplazar completamente la comunicación directa, sino disponer de una fuente central de información que permita a cada actor consultar los datos que le corresponden según su rol.

\subsubsection{Acceso digital}
\label{sec:acceso-digital}

La disponibilidad de conectividad constituye un factor crítico a considerar, especialmente en una plataforma que integra funciones de localización y seguimiento. Según datos de la Unión Internacional de Telecomunicaciones, en Nicaragua existe una cobertura significativa de redes móviles; para el año 2022, se reportaba una cobertura de al menos 3G del 85~\% y de 4G del 73~\%, aunque con disparidades entre tecnologías \parencite{ituDigitalNicaragua}. No obstante, la existencia de cobertura no garantiza una calidad de servicio constante. TELCOR reportó para 2024 una cobertura móvil nacional del 87,7~\%, señalando que varios indicadores de calidad requieren mejoras \parencite{telcor2024calidad}. Ante este panorama, TravelDLS debe prever que la conexión puede variar a lo largo de un recorrido. Por ello, la plataforma se concibe como una aplicación web adaptable a diferentes tamaños de pantalla, diseñada para gestionar adecuadamente los retrasos o interrupciones en las actualizaciones de ubicación. Asimismo, la visualización de la fecha y hora de la última posición recibida resulta fundamental para evitar que una ausencia temporal de datos sea interpretada erróneamente como una detención del vehículo.

\subsubsection{Transparencia y confianza}
\label{sec:transparencia-confianza}

La confianza en un servicio de transporte también se relaciona con la calidad y disponibilidad de la información durante su ejecución. \textcite{nie2026visibility} distinguen entre visibilidad, trazabilidad y transparencia: la visibilidad se relaciona con disponer de información actual sobre una operación, la trazabilidad con conservar información que permita reconstruir lo ocurrido y la transparencia con proporcionar información adecuada a los actores involucrados.

Esta relación entre información y confianza también se observa en el estudio de \textcite{brookbanks2022impact}, quienes encontraron que una plataforma digital basada en datos compartidos y confiables puede reducir la duplicación de información, mejorar la visibilidad de la cadena de suministro y apoyar la construcción de confianza entre sus participantes. Los autores también aclaran que la tecnología no elimina la necesidad de establecer relaciones de confianza entre los actores involucrados.

En TravelDLS, estos principios se reflejan en la consulta del estado y ubicación de los servicios. La posición más reciente del vehículo aporta visibilidad, mientras que el registro de estados y ubicaciones contribuye a la trazabilidad de la operación. A su vez, el cliente podrá consultar la información correspondiente al servicio contratado, mientras que los mecanismos de autenticación y roles limitarán el acceso según las responsabilidades de cada usuario. De esta manera, la transparencia no implica mostrar toda la información disponible, sino proporcionar datos relevantes y confiables sin comprometer la privacidad de los participantes.

En conjunto, estos aspectos muestran que la incorporación de TravelDLS no debe analizarse únicamente desde su funcionamiento técnico, sino también desde la manera en que sus posibles usuarios podrían relacionarse con la plataforma durante sus actividades cotidianas. La adopción progresiva de herramientas digitales, la organización de los procesos de trabajo, la coordinación entre empresa, conductor y cliente, las condiciones variables de conectividad y la necesidad de disponer de información confiable influyen directamente en las decisiones de diseño del sistema. Por esta razón, funcionalidades como la centralización de la información, la organización por roles, una interfaz adaptable y el seguimiento de las operaciones no responden solamente a requerimientos funcionales, sino también a las condiciones identificadas en la literatura sobre el uso de tecnologías digitales. Considerados estos factores, resulta necesario profundizar en los conceptos y fundamentos que explican los componentes principales de TravelDLS, los cuales se desarrollan a continuación en el Marco Teórico.
% Capítulo 2: Marco Teórico
% Propósito: Presentar las bases teóricas, conceptos clave, antecedentes de investigación y el estado del arte que fundamentan el proyecto.

\section{Marco Teórico}
\label{sec:marco-teorico}

El Marco Teórico reúne los antecedentes, conceptos y fundamentos necesarios para comprender los principales elementos que forman parte de TravelDLS. En este apartado se analizan investigaciones relacionadas con la propuesta, se definen los términos utilizados dentro del proyecto y se presentan las bases teóricas que sustentan la gestión del transporte, el seguimiento y la planificación de rutas.

\subsection{Antecedentes del Proyecto}
\label{sec:antecedentes-proyecto}

Los trabajos revisados en los \textbf{\hyperref[sec:antecedentes]{Antecedentes}} ofrecen referentes para la gestión del transporte, el seguimiento de la carga y la planificación de rutas. A partir de sus aportes y de otras investigaciones relacionadas, se examinan las diferencias de TravelDLS y su contribución propuesta al servicio de carga pesada en Nicaragua.

En la gestión del transporte, el \textbf{\hyperref[ant:aragon]{sistema de administración de fletes de Master Logistics}} organiza la información operativa de una empresa \parencite[pp.~2, 9]{aragon2018implementacion}. TravelDLS retoma esta función para atender a distintas transportistas, manteniendo separados sus registros y vinculándolos con la contratación y el seguimiento de los viajes. Con ello, cada empresa podrá compartir con el cliente el estado de su servicio.

Por su parte, ROUTFALL vincula los registros de clientes, pedidos y costos con la elaboración de rutas para empresas de La Unión, El Salvador \parencite[pp.~52--53, 59--60]{flores2017rutas}. Mientras este trabajo se centra en organizar la distribución, TravelDLS plantea un proceso que comprende la publicación, la contratación y el seguimiento de servicios de carga pesada. La planificación del recorrido formará parte de ese proceso para facilitar la coordinación entre empresa, conductor y cliente.

En cuanto al seguimiento, el \textbf{\hyperref[ant:garcia]{sistema de localización de unidades en León}} permite consultar la ubicación del transporte de pasajeros \parencite[pp.~3, 7]{garcia2022desarrollo}. TravelDLS relacionará esa ubicación con el estado del viaje contratado para que el cliente pueda consultar ambos aspectos del traslado de su carga. El \textbf{\hyperref[ant:pauzuoliene]{estudio de adopción tecnológica en empresas logísticas}} respalda la pertinencia de estas funciones, aunque no demuestra su aceptación en Nicaragua \parencite[pp.~988--989]{pauzuoliene2024smart}.

Los estudios de \textbf{\hyperref[ant:aguilar]{distribución en Tegucigalpa}} y de \textbf{\hyperref[ant:hernandez]{software para Transporte Gonzales}} evalúan herramientas de planificación de rutas en escenarios concretos \parencites[p.~56]{aguilar2022optimizacion}[p.~111]{hernandez2024evaluacion}. TravelDLS seguirá ese criterio al comparar Vecino Más Cercano y 2-opt con Google OR-Tools. Además de evaluar los métodos, integrará esta última herramienta en la gestión de los viajes contratados, sin asumir que alcanzará los ahorros de esos estudios.

En esta misma línea, \textcite[pp.~15--17, 28--31]{martinez2020ruteo} abordan la distribución de bebidas en Quito como un problema de ruteo con capacidad y ventanas de tiempo (CVRPTW), mediante agrupamiento k-medoids y un algoritmo genético. En su evaluación obtuvieron una menor distancia respecto de las estimaciones de la operación estudiada \parencite[p.~37]{martinez2020ruteo}. TravelDLS se diferencia en el problema y los métodos seleccionados: se limita al TSP y a las alternativas antes mencionadas. Por tanto, este antecedente orienta la evaluación de recorridos, pero no constituye una referencia de rendimiento directamente comparable.

Como referencia complementaria, \textcite[pp.~20--24]{alvarez2016plataforma} modelan la distribución de lechuga mediante la asignación de vehículos y la ubicación de centros logísticos. Su plataforma es una red física, mientras que TravelDLS es una aplicación web para gestionar servicios y viajes. Las limitaciones de los datos y del supuesto de transportar un solo producto, reconocidas por los autores, muestran la necesidad de interpretar los resultados según las condiciones de cada modelo.

Entre las soluciones más cercanas a la propuesta, la \textbf{\hyperref[ant:sultan]{plataforma de Sultan y colaboradores}} integra reservas, seguimiento y gestión logística, e incluye conductores que eligen pedidos disponibles para reparto \parencite[pp.~1--4]{sultan2022real}. TravelDLS plantea un esquema distinto para la carga pesada: cada transportista asignará sus vehículos y conductores a los viajes contratados. La diferencia radica en mantener las asignaciones bajo la administración de cada empresa, con acceso del cliente a la información de su servicio, y no en integrar funciones que el antecedente ya reúne.

Asimismo, \textcite{lasluisa2026diseno} utilizan el modelo SaaS multitenant para gestionar flotas y compañías en un sistema de rastreo satelital. TravelDLS comparte este enfoque y lo orienta a conectar la publicación y contratación del servicio con el seguimiento del viaje. El aporte esperado está en la continuidad entre la relación comercial y la operación de transporte, no en el uso de una arquitectura multitenant. El alcance consultado de este antecedente no permite descartar que incluya funciones equivalentes.


Además de las investigaciones relacionadas con la gestión del transporte, el seguimiento de unidades y la planificación de rutas, se consideran estudios que han utilizado tecnologías para desarrollar aplicaciones web con diferentes funcionalidades. Estos trabajos permiten conocer las razones por las que sus autores seleccionaron determinadas herramientas y los aportes que pueden retomarse para el desarrollo de TravelDLS.

En relación con el desarrollo de la plataforma, \textcite{nguyen2022building} utilizó Next.js en una aplicación de comercio electrónico por sus capacidades de renderizado del lado del servidor y generación de páginas estáticas. Asimismo, \textcite{kolding2021developing} incorporaron Next.js en BikeViz para organizar y presentar información geográfica mediante una aplicación web. Ambos proyectos aportan referentes para el desarrollo de la interfaz de TravelDLS, que permitirá publicar servicios de transporte y consultar información relacionada con
las operaciones.

Para el desarrollo del backend, \textcite{farinha2025arkivo} seleccionó AdonisJS en la construcción de un sistema de preservación de arte digital debido a su arquitectura MVC, su compatibilidad con TypeScript y las convenciones que facilitan la organización del código. Esta experiencia constituye un referente para la organización de los módulos de TravelDLS. En cuanto al almacenamiento de información,
\textcite{salunke2024performance} compararon PostgreSQL y MySQL, obteniendo resultados favorables para PostgreSQL en las consultas y operaciones concurrentes evaluadas. Su
investigación aporta criterios relacionados con el rendimiento de las bases de datos y la gestión de los registros de la plataforma.

En cuanto a la representación geográfica, \textcite{kolding2021developing} también seleccionaron Mapbox GL JS para BikeViz debido a sus posibilidades de personalización, la documentación disponible y la facilidad para incorporar funciones a los mapas. Mediante esta herramienta desarrollaron una interfaz que permitía explorar información geográfica a través de diferentes capas. Este antecedente se relaciona con TravelDLS, donde Mapbox se utilizará para representar las ubicaciones de los vehículos y los recorridos asociados con los viajes.

En conjunto, los antecedentes revisados permiten identificar aportes relacionados con la administración del transporte, la contratación de servicios, la asignación de recursos, el seguimiento de viajes y la planificación de recorridos. Asimismo, las investigaciones tecnológicas presentan experiencias sobre la selección de herramientas para el desarrollo de aplicaciones web y la gestión de información. TravelDLS retoma estos aportes y los adapta al transporte terrestre de carga pesada en Nicaragua, mediante una plataforma que relacionará la solicitud del servicio con la administración y el seguimiento de cada viaje hasta su finalización.


\subsection{Definición de Términos}

Para comprender la propuesta de TravelDLS es necesario establecer algunos conceptos que se encuentran directamente relacionados con las operaciones que se pretenden gestionar dentro de la plataforma.

La \hypertarget{concepto-gestion-logistica}{\textbf{gestión logística}} constituye uno de los conceptos centrales del proyecto, ya que una operación de transporte no depende únicamente del desplazamiento de una carga. También requiere organizar información sobre vehículos, conductores, recorridos y estados del servicio. Dentro de TravelDLS, este concepto se entenderá como la coordinación de la información y de las actividades necesarias para desarrollar un servicio de transporte desde su asignación hasta el seguimiento del viaje. En sistemas logísticos desarrollados anteriormente, la integración de clientes, conductores, reservas, envíos y seguimiento ha permitido mantener diferentes componentes de la operación dentro de un mismo entorno \parencite[pp. 1-4]{sultan2022real}. Esta idea se relaciona con TravelDLS, donde se busca que los datos del viaje no se manejen de forma aislada, sino como partes de una misma operación.


En este proceso, la \hypertarget{concepto-orden}{\textbf{orden de transporte o servicio de flete}} representa la solicitud de un servicio para trasladar una carga entre los lugares establecidos por el cliente y la empresa transportista. En este registro se relacionan los datos generales de la carga, su origen y destino, la empresa que prestará el servicio, el precio acordado y el estado de la operación. De esta manera, una misma solicitud puede consultarse desde las distintas funciones del sistema, sin manejar por separado la información de contratación, asignación y seguimiento. La necesidad de organizar los procesos de fletes mediante sistemas de información también ha sido abordada en Nicaragua por \textcite{aragon2018implementacion}.

De igual manera, el \hypertarget{concepto-estados}{\textbf{ciclo de atención y los estados de una orden}} permiten representar las distintas etapas por las que pasa una solicitud desde su registro hasta la conclusión del servicio. Este proceso puede incluir la cotización, la aceptación o el rechazo del precio, la asignación de los recursos necesarios y el desarrollo del traslado. En TravelDLS, los estados ayudarán a identificar la situación actual de cada operación y a diferenciar el avance del servicio de los cambios asociados al pago, ya que una orden entregada y una transacción pagada describen situaciones distintas. La organización de pedidos en sistemas de gestión del transporte se describe en \textcite{sepulveda2023propuesta}; por ello, los estados específicos deberán ajustarse a las reglas del negocio que finalmente se aprueben.

Para ejecutar la orden, la \hypertarget{concepto-asignacion}{\textbf{asignación de recursos de transporte}} consiste en vincular la solicitud con los medios que participarán en su ejecución, principalmente el vehículo y el conductor. Esta relación permite que la empresa identifique quién atenderá la orden y a qué unidad corresponde el traslado, además de consultar la información del recorrido y de la carga desde el mismo servicio. La coordinación de flotas y conductores en el transporte de carga pesada se aborda en \textcite{laura2024tecnologias}. En el diseño actual de TravelDLS se contempla registrar categorías de camión, conductores y vehículos; sin embargo, las reglas específicas sobre disponibilidad, capacidad y compatibilidad deberán establecerse conforme al alcance real del proyecto.

A medida que un viaje se encuentra en desarrollo, cobra importancia conocer qué ocurre con el servicio y conservar información sobre su avance. De esta necesidad surgen los conceptos de \hypertarget{concepto-trazabilidad}{\textbf{trazabilidad}} y \hypertarget{concepto-seguimiento-tiempo-real}{\textbf{seguimiento en tiempo real}}. La trazabilidad permite mantener información relacionada con los estados y acontecimientos de un servicio durante su recorrido, mientras que el seguimiento en tiempo real se concentra en ofrecer la ubicación o la información más reciente disponible. \textcite{pauzuoliene2024smart} señalan que las tecnologías de seguimiento permiten conocer la ubicación y el estado de la carga, además de facilitar la estimación de los tiempos de entrega (p. 988). En TravelDLS, estas funciones estarán relacionadas tanto con las necesidades de control de la empresa transportista como con la información que necesita consultar el cliente durante el traslado.


Junto con la información más reciente del viaje, el \hypertarget{concepto-historial}{\textbf{historial de seguimiento}} permite conservar registros sucesivos para reconstruir la secuencia de ubicaciones o acontecimientos de un envío. Mientras la ubicación actual muestra la posición reportada durante el recorrido, el historial reúne los datos almacenados a lo largo del tiempo. En TravelDLS, esta distinción permitirá que el conductor reporte información vinculada a una orden y que el cliente consulte el progreso de su carga. La frecuencia y oportunidad de las actualizaciones dependerán de la conexión a Internet, la disponibilidad del dispositivo y el mecanismo de registro empleado. Como antecedente nacional de localización de unidades de transporte, se considera el trabajo de \textcite{garcia2022desarrollo}.

El seguimiento de una unidad requiere conocer su posición durante el recorrido, por lo que la \hypertarget{concepto-geolocalizacion}{\textbf{geolocalización}} constituye otro concepto relacionado directamente con el proyecto. Esta permite obtener y representar la ubicación geográfica de un vehículo mediante tecnologías de posicionamiento y servicios de mapas. En Nicaragua, \textcite{garcia2022desarrollo} incorporaron servicios de Google Maps en una aplicación destinada a mostrar en tiempo real la ubicación de unidades de transporte (p. 7). Aunque dicho proyecto corresponde al transporte público, demuestra la aplicación de tecnologías de localización dentro de un sistema de transporte desarrollado en el contexto nacional. En TravelDLS, la geolocalización estará vinculada al seguimiento del viaje y a la consulta de la ubicación de la carga.

Además de conocer dónde se encuentra un vehículo, resulta necesario determinar el recorrido que debe seguir. En este contexto aparecen la \hypertarget{concepto-planificacion-rutas}{\mbox{\textbf{planificación de rutas}}} y la \hypertarget{concepto-optimizacion-rutas}{\mbox{\textbf{optimización de rutas}}}. La planificación consiste en definir y ordenar los puntos que deben visitarse durante un recorrido, mientras que la optimización busca comparar distintas alternativas con el propósito de encontrar recorridos que presenten mejores resultados según criterios como distancia, tiempo o costos. Los antecedentes regionales muestran que los resultados obtenidos mediante herramientas de optimización pueden variar según las condiciones del problema y la solución utilizada. \textcite{aguilar2022optimizacion} reportó mejoras en el escenario analizado mediante la reducción de la distancia recorrida (p. 56), mientras que \textcite{hernandez2024evaluacion} obtuvieron resultados diferentes al evaluar alternativas de software para planificación de rutas (p. 111). Estas diferencias justifican que en TravelDLS no se asuma desde el inicio que un método será superior, sino que se realice una comparación entre las alternativas definidas para el proyecto.


Para comparar recorridos entre varios destinos, se puede utilizar una \hypertarget{concepto-matriz}{\textbf{matriz de distancias}}, una estructura organizada en filas y columnas en la que cada posición representa la distancia entre un punto de origen y otro de destino. Esta representación permite que los métodos de optimización utilicen los mismos datos para calcular y comparar el costo total de distintas rutas. \textcite{tuo2024matrix} emplearon una matriz de distancias dentro de un modelo de distribución logística que incorpora información sobre las condiciones de las vías. Asimismo, cuando el criterio de análisis es la duración del trayecto, la matriz puede representar tiempos de desplazamiento; \textcite{lombard2018matrix} estudiaron un problema de ruteo con tiempos de viaje variables y utilizaron una matriz de distribuciones de probabilidad en lugar de una matriz tradicional de distancias. En TravelDLS, la matriz de distancias o tiempos se contempla como insumo para comparar recorridos con múltiples destinos mediante Vecino Más Cercano, 2-opt y OR-Tools. La forma de obtener estos valores deberá definirse durante la implementación según los datos cartográficos y las condiciones consideradas en el prototipo.

La optimización de recorridos conduce al \hypertarget{concepto-tsp}{\textbf{Problema del Agente Viajero}}, conocido como TSP por sus siglas en inglés. Este problema plantea la necesidad de visitar un conjunto de puntos y regresar al lugar de origen procurando reducir la distancia total recorrida. Dentro de las técnicas que pueden utilizarse para construir una solución se encuentra el \hypertarget{concepto-vecino-mas-cercano}{\textbf{Vecino Más Cercano}}, que selecciona progresivamente el siguiente destino tomando como referencia la menor distancia disponible desde el punto actual. \textcite{contreras2021aplicacion} describen la utilización de este método en situaciones relacionadas con transporte y generación de rutas (p. 11). Para TravelDLS, estos conceptos resultan especialmente importantes debido a que el proyecto contempla el análisis de viajes con múltiples destinos y la comparación de diferentes métodos de optimización.


Otra técnica relacionada con la construcción de recorridos es el \hypertarget{concepto-dos-opt}{\textbf{algoritmo 2-opt}}, un procedimiento de búsqueda local que intenta mejorar una ruta ya construida. Para ello, selecciona dos conexiones del recorrido, las sustituye por una combinación diferente y reorganiza el tramo correspondiente para comprobar si disminuye la distancia total. Cuando encuentra una alternativa de menor costo, conserva el cambio y continúa buscando mejoras hasta alcanzar el criterio de finalización establecido. \textcite{uddin2023twoopt} describen este procedimiento en una propuesta aplicada al Problema del Agente Viajero, mientras que \textcite{dacosta2021twoopt} estudian operadores 2-opt como parte de estrategias de búsqueda local para problemas de rutas. En TravelDLS, 2-opt se considera un procedimiento complementario al Vecino Más Cercano: este último construye una ruta inicial mediante elecciones sucesivas y 2-opt puede utilizarla como punto de partida para mejorar el orden de visita. Los resultados de esta combinación podrán compararse con los obtenidos mediante OR-Tools usando los mismos puntos y criterios de evaluación. No obstante, 2-opt no garantiza encontrar siempre la ruta de menor distancia posible, por lo que su utilidad deberá valorarse mediante las pruebas del proyecto \parencite{dacosta2021twoopt,uddin2023twoopt}.

Estos conceptos permiten comprender los principales procesos que TravelDLS busca integrar. Sin embargo, para explicar de manera más amplia cómo estos elementos pueden organizarse dentro de una plataforma tecnológica, es necesario relacionarlos con enfoques teóricos sobre gestión de transporte, transporte de carga pesada y prestación de software.

\subsection{Bases Teóricas}

Las bases teóricas de TravelDLS permiten explicar los enfoques que sustentan tanto la gestión de las operaciones de transporte como la forma en que la plataforma será ofrecida a diferentes empresas. Uno de los principales referentes son los \hypertarget{concepto-tms}{}Sistemas de Gestión de Transporte (TMS), los cuales integran actividades como la planificación, ejecución y seguimiento de los envíos. \textcite[pp.~10--11]{sepulveda2023propuesta} señala que este tipo de sistemas puede incluir funciones como la gestión de pedidos, planificación de rutas, asignación de cargas, seguimiento de vehículos y facturación, además de facilitar el análisis de indicadores relacionados con costos, tiempos de entrega y utilización de recursos. Estas características guardan relación con TravelDLS, ya que la plataforma busca mantener conectada la información de viajes, vehículos, conductores, cargas y rutas dentro de un mismo entorno.

Esta necesidad de integración se vuelve más importante al considerar las características propias del \hypertarget{concepto-carga-pesada}{}transporte terrestre de carga pesada, donde las operaciones requieren coordinar vehículos, conductores, recorridos y tiempos de entrega. \textcite[p.~67]{laura2024tecnologias} explica que este tipo de transporte presenta desafíos relacionados con la gestión de flotas, la planificación de rutas y la coordinación de conductores, además de restricciones como límites de carga y horarios de conducción. En este contexto, la tecnología sirve como apoyo para centralizar información y automatizar ciertos procesos, lo cual se relaciona directamente con la propuesta de TravelDLS de reunir varias funciones operativas dentro de una misma plataforma.

La forma en que TravelDLS será ofrecido también encuentra sustento en el modelo \hypertarget{concepto-saas}{}Software como Servicio (SaaS). Bajo este enfoque, la aplicación se encuentra disponible a través de Internet y el proveedor se encarga del mantenimiento y las actualizaciones, evitando que cada organización deba administrar una instalación independiente. \textcite[p.~33]{lasluisa2026diseno} explican que este modelo permite ofrecer el software a múltiples usuarios mediante una infraestructura común. Esta característica resulta especialmente adecuada para TravelDLS, ya que la plataforma está pensada para ser utilizada por distintas empresas transportistas y no por una sola organización.

Relacionado con lo anterior, la \hypertarget{concepto-multitenant}{}arquitectura multitenant permite que varias organizaciones utilicen una misma aplicación manteniendo separada la información correspondiente a cada una. \textcite[pp.~34--36]{lasluisa2026diseno} señalan que este enfoque facilita el uso compartido de recursos y el mantenimiento del sistema, aunque también exige mecanismos de seguridad y control de acceso que eviten que un tenant pueda consultar información que no le corresponde. En TravelDLS, esta separación será necesaria para que cada empresa administre sus propios vehículos, conductores, viajes y cargas, mientras que los clientes solo puedan acceder a la información relacionada con los servicios que han contratado.

Esta característica también requiere considerar la \hypertarget{concepto-separacion}{}separación de información operativa, debido a que las empresas transportistas utilizarán una plataforma compartida, pero cada una deberá mantener diferenciados sus registros. En TravelDLS, esta organización se refleja en las relaciones existentes entre empresas, vehículos, conductores y órdenes, permitiendo identificar a qué organización pertenece cada recurso y qué usuarios están autorizados para consultar o modificar sus datos. La importancia de aplicar mecanismos de control de acceso dentro de entornos multitenant se relaciona con lo expuesto por \textcite[pp.~34--36]{lasluisa2026diseno}, quienes destacan la necesidad de proteger la información de las diferentes organizaciones que utilizan una infraestructura común. De esta manera, TravelDLS buscará combinar el uso compartido de la plataforma con una administración diferenciada de los datos operativos de cada empresa.

Además de los enfoques relacionados con la administración y prestación del software, TravelDLS incorpora fundamentos de optimización matemática para apoyar la planificación de recorridos. En este ámbito se considera \hypertarget{concepto-ortools}{}Google OR-Tools, un conjunto de herramientas de software desarrollado para modelar y resolver diferentes problemas de optimización, incluyendo aquellos relacionados con la planificación de rutas. A diferencia del Vecino Más Cercano y 2-opt, que son procedimientos específicos para construir o mejorar recorridos, OR-Tools proporciona componentes que permiten representar el problema, establecer restricciones y ejecutar distintas estrategias de búsqueda. \textcite{muriyatmoko2024ortools} utilizaron esta herramienta con Python para comparar heurísticas y metaheurísticas aplicadas a un problema de ruteo de vehículos con restricciones de capacidad, mientras que \textcite{salsabila2023ortools} la emplearon en un escenario de distribución logística para analizar estrategias como Guided Local Search y Simulated Annealing.

La documentación oficial de Google señala que OR-Tools permite trabajar con diferentes configuraciones de problemas de ruteo y aplicar procedimientos de búsqueda para obtener soluciones de acuerdo con las restricciones establecidas, aunque estas no siempre garantizan alcanzar el resultado globalmente óptimo \parencite{googleortoolsRouting}. Para TravelDLS, la incorporación de esta herramienta está relacionada con el objetivo de desarrollar un módulo que apoye la planificación de rutas con múltiples destinos, utilizando datos de desplazamiento organizados mediante matrices de distancias o tiempos. Su aplicación permitirá comparar los recorridos generados mediante Vecino Más Cercano y mejorados con 2-opt frente a los resultados obtenidos con OR-Tools, manteniendo criterios de evaluación equivalentes. Esta comparación permitirá analizar el comportamiento de las alternativas consideradas antes de valorar su integración y desempeño dentro del prototipo.

En conjunto, estas bases teóricas permiten entender que TravelDLS no se limita a ser una aplicación para registrar viajes, sino que se apoya en principios de gestión de transporte, centralización de procesos y prestación de software mediante una plataforma compartida. Los TMS aportan el enfoque de organización y seguimiento, el transporte de carga pesada define el contexto operativo y el modelo SaaS junto con la arquitectura multitenant explican cómo el sistema puede ser utilizado por diferentes empresas dentro de una misma solución tecnológica. A estos fundamentos se incorpora la optimización de rutas, que permite analizar recorridos con múltiples destinos mediante procedimientos matemáticos y herramientas especializadas.

Con estos fundamentos se establece la base conceptual que orienta la propuesta de TravelDLS. A partir de ella, el siguiente paso consiste en definir el procedimiento metodológico utilizado para seleccionar y analizar la información que sustenta el proyecto, así como organizar las actividades necesarias para continuar con su desarrollo.

% Capítulo 3: Metodología
% Propósito: Detallar el enfoque metodológico, diseño de investigación,
% técnicas e instrumentos de recolección de información y la metodología
% de desarrollo de software seleccionada.

\section{Metodología}
\label{sec:metodologia}

\subsubsection{Enfoque metodológico y tipología del artefacto}
\label{sec:metodo-enfoque}

Para orientar el desarrollo de TravelDLS se adopta como enfoque general la Metodología de Investigación en Ciencia del Diseño, conocida por sus siglas en inglés como DSRM (\textit{Design Science Research Methodology}), propuesta por \textcite{peffers2007dsrm}. Este enfoque resulta pertinente para el proyecto porque la investigación no se limita a describir las dificultades relacionadas con la administración y el seguimiento del transporte terrestre de carga pesada, sino que también busca desarrollar una solución informática que responda a las necesidades identificadas.

La metodología DSRM permite relacionar la identificación de un problema con la definición de los objetivos de la solución, su diseño, construcción y posterior comprobación. Dentro de TravelDLS, esta relación es importante porque permite mantener una conexión entre la problemática que motiva el proyecto y las funcionalidades que se incorporarán a la plataforma. De esta manera, las decisiones técnicas no se consideran actividades aisladas, sino parte de un proceso orientado a responder a las necesidades de gestión, localización, seguimiento y planificación de rutas.

En el ámbito de los sistemas de información, \textcite{hevner2004design} explican que la investigación basada en Ciencia del Diseño puede generar diferentes tipos de artefactos, entre ellos constructos, modelos, métodos e instanciaciones. Estos elementos representan distintas formas de conocimiento que contribuyen a comprender un problema y desarrollar soluciones tecnológicas. Para TravelDLS, el artefacto principal previsto corresponde a una instanciación, representada por la plataforma web que integrará las funciones de administración de empresas transportistas, conductores, vehículos, cargas, órdenes, pagos y seguimiento de los servicios.

La construcción de esta instanciación se apoya en tres elementos complementarios. Los constructos proporcionan los conceptos y el vocabulario necesarios para comprender el dominio de aplicación. Los modelos permiten representar la estructura de la información y el comportamiento esperado del sistema. Por su parte, los métodos comprenden los procedimientos que se utilizarán para organizar la investigación, desarrollar las funcionalidades y resolver problemas específicos, como la planificación de recorridos con múltiples destinos.

Esta organización permite que los fundamentos desarrollados en el marco teórico tengan una aplicación dentro del marco metodológico, evitando repetir definiciones y manteniendo una relación entre los conceptos estudiados, los modelos elaborados y las decisiones que orientarán el desarrollo de TravelDLS.

\subsubsection{Fase 1. Identificación del Problema y Motivación}
\label{sec:metodo-fase1}

La primera fase de la metodología DSRM se orienta a identificar el problema que motiva la investigación y fundamentar la necesidad de desarrollar una solución tecnológica \parencite{peffers2007dsrm}. En TravelDLS, esta etapa parte del análisis de las actividades que intervienen en la gestión del transporte terrestre de carga pesada y de las dificultades que surgen cuando la información necesaria para desarrollar un servicio no se encuentra debidamente relacionada.

El transporte de carga pesada comprende diferentes actividades que deben coordinarse para completar una operación. Además del traslado físico de la mercancía, es necesario administrar solicitudes, vehículos, conductores, cargas, recorridos y comunicaciones entre los participantes del servicio. Estos elementos mantienen una relación directa, debido a que una operación requiere identificar quién contrata el transporte, qué empresa prestará el servicio, cuáles son los recursos asignados y en qué situación se encuentra el traslado.

Cuando esta información se administra mediante registros separados, hojas de cálculo, llamadas o mensajes individuales, resulta más difícil mantener una visión organizada del servicio. Para consultar el avance de una operación es necesario reunir datos provenientes de diferentes medios, lo que dificulta identificar con rapidez su estado y los recursos asociados. Esta situación evidencia la importancia de establecer mecanismos que permitan centralizar la información y mantener una relación entre las diferentes actividades operativas.

En Nicaragua existen investigaciones que han abordado necesidades relacionadas con la administración del transporte mediante herramientas informáticas. \textcite{aragon2018implementacion} desarrollaron un sistema de información para la administración de fletes terrestres, orientado a mejorar la organización de los registros y procesos correspondientes a una empresa de transporte. Este antecedente permite reconocer la importancia que tiene la administración de la información dentro de las operaciones logísticas.

De manera complementaria, \textcite{garcia2022desarrollo} abordaron el seguimiento y la localización de unidades de transporte público en la ciudad de León mediante una aplicación móvil. Aunque este trabajo corresponde a un contexto diferente al transporte de carga pesada, permite reconocer la utilidad de las tecnologías de localización para consultar información sobre el desplazamiento de las unidades.

Estos antecedentes contribuyen a comprender dos necesidades relacionadas con TravelDLS. La primera corresponde a la organización de los registros y procesos administrativos de las empresas transportistas. La segunda se relaciona con la posibilidad de consultar información sobre la ubicación y el desplazamiento de los vehículos durante la prestación de un servicio.

Desde la perspectiva administrativa, una empresa dedicada al transporte de carga pesada necesita mantener relacionados los datos de los recursos que utiliza y las solicitudes que recibe. No resulta suficiente disponer de un registro de conductores y vehículos si esa información no puede relacionarse con las órdenes de transporte que se encuentran pendientes o en ejecución. Por esta razón, la propuesta contempla que las operaciones se organicen mediante una \textbf{\hyperlink{concepto-orden}{orden de transporte}}, que permita vincular los datos del servicio con la empresa, el cliente y los recursos correspondientes.

Asimismo, las solicitudes requieren pasar por diferentes etapas, desde su creación hasta la conclusión del servicio. Esto introduce la necesidad de establecer un \textbf{\hyperlink{concepto-estados}{ciclo de atención}} \textbf{\hyperlink{concepto-estados}{y estados de una orden}}, mediante el cual se identifique cuándo una solicitud se encuentra pendiente, cuándo ha recibido un precio y cuándo ha sido aceptada o rechazada por el cliente. La definición de estos estados permitirá representar el avance administrativo y operativo de una solicitud, manteniendo una relación entre las acciones realizadas y la situación actual del servicio.

Otra dificultad relacionada con la administración del transporte corresponde a la \textbf{\hyperlink{concepto-asignacion}{asignación de recursos}}. Cada solicitud requiere identificar los medios que participarán en su atención, principalmente el conductor y el vehículo. Cuando esta relación no se mantiene dentro de un mismo registro, la consulta de la operación depende de información adicional que puede encontrarse en medios diferentes. TravelDLS busca organizar estos datos para que la empresa pueda identificar los recursos asociados a las órdenes que administra.

Desde la perspectiva del cliente, una de las principales necesidades consiste en disponer de información sobre el servicio contratado. Después de aceptar una solicitud de transporte, resulta importante conocer si la carga ha comenzado a trasladarse, cuál es el estado del servicio y dónde se encuentra la unidad encargada del recorrido.

Cuando esta información no puede consultarse de manera directa, el cliente depende de la comunicación individual con la empresa o el conductor para conocer el avance de su carga. Esto puede generar incertidumbre sobre la situación del traslado, especialmente cuando no existen mecanismos que permitan consultar la ubicación reportada y los cambios de estado de la operación.

La incorporación de funciones de \textbf{\hyperlink{concepto-trazabilidad}{trazabilidad}} y \textbf{\hyperlink{concepto-seguimiento-tiempo-real}{seguimiento en tiempo real}} responde a esta necesidad. Mediante estas funcionalidades se pretende mantener vinculados los registros de ubicación y los estados de la operación con la orden correspondiente, permitiendo que los usuarios autorizados consulten información relacionada con el desarrollo del servicio.

La \textbf{\hyperlink{concepto-geolocalizacion}{geolocalización}} constituye un elemento necesario para representar la posición de la unidad durante el traslado. Sin embargo, conocer únicamente la ubicación más reciente no permite reconstruir los desplazamientos realizados anteriormente. Por esta razón, también se contempla el \textbf{\hyperlink{concepto-historial}{historial de seguimiento}}, que permitirá conservar registros sucesivos asociados a una misma operación.

Esta diferencia resulta importante porque el seguimiento de un servicio no depende solamente de mostrar un punto sobre un mapa. También requiere relacionar la ubicación reportada con el conductor, la orden correspondiente y el estado del traslado. De esta manera, la información consultada podrá interpretarse dentro del contexto de la operación que se está desarrollando.

No obstante, el seguimiento estará condicionado por la disponibilidad de conexión a Internet, los permisos de localización y las características de los dispositivos utilizados. Estas condiciones deben considerarse al establecer las funcionalidades y limitaciones del proyecto, evitando afirmar que la plataforma garantizará una actualización continua de ubicaciones bajo cualquier circunstancia.

Además de la administración y el seguimiento, otra necesidad identificada corresponde a la organización de recorridos con múltiples destinos. Una empresa transportista puede requerir visitar diferentes puntos durante un mismo servicio, por lo que resulta pertinente analizar el orden en que se realizará el desplazamiento.

La \textbf{\hyperlink{concepto-planificacion-rutas}{planificación de rutas}} organiza los destinos. La \mbox{\textbf{\hyperlink{concepto-optimizacion-rutas}{optimización de rutas}}} permite comparar recorridos según criterios como la distancia o el tiempo. Esto resulta relevante debido a que el orden de visita influye en la distancia total que puede representar una operación.


En el contexto nacional, \textcite{aragon2018implementacion}
abordaron la administración de fletes terrestres mediante
un sistema de información, mientras que
\textcite{garcia2022desarrollo} estudiaron el seguimiento
y localización de unidades de transporte. Aunque estas
investigaciones no se centran directamente en la
optimización de rutas de carga pesada, permiten
reconocer la aplicación de herramientas tecnológicas
en el sector del transporte nicaragüense.

Por otra parte, los antecedentes regionales de
\textcite{aguilar2022optimizacion} y
\textcite{hernandez2024evaluacion} abordan la
planificación y optimización de rutas mediante
diferentes alternativas. Estas investigaciones
constituyen una referencia para el componente de
planificación de recorridos que se contempla
desarrollar en TravelDLS.


Estas investigaciones permiten fundamentar la pertinencia de estudiar alternativas de optimización, aunque los resultados obtenidos corresponden a los escenarios particulares de cada trabajo y no pueden considerarse mejoras garantizadas para TravelDLS.

Por esta razón, el proyecto contempla analizar el \textbf{\hyperlink{concepto-tsp}{Problema del Agente Viajero}} y comparar procedimientos como \textbf{\hyperlink{concepto-vecino-mas-cercano}{Vecino Más Cercano}} y \textbf{\hyperlink{concepto-dos-opt}{2-opt}} con las soluciones generadas mediante \textbf{\hyperlink{concepto-ortools}{Google OR-Tools}}. El propósito de esta comparación será conocer el comportamiento de las alternativas consideradas y establecer una base técnica para el módulo de planificación de rutas del prototipo.

La problemática también se relaciona con la presencia tecnológica de las empresas transportistas. TravelDLS contempla reunir diferentes organizaciones dentro de una misma plataforma, permitiendo que puedan presentar sus servicios a clientes interesados en contratar transporte de carga pesada. Esta característica busca conectar la administración interna con las actividades de búsqueda y contratación, ofreciendo un espacio común para empresas y clientes.

La posibilidad de mejorar la visibilidad de las empresas constituye un beneficio esperado de la propuesta. Sin embargo, no debe presentarse como un resultado demostrado, debido a que todavía sería necesario evaluar el uso de la plataforma y su aceptación dentro del entorno operativo correspondiente.

Como respuesta a las necesidades identificadas, TravelDLS propone desarrollar una plataforma bajo el modelo \textbf{\hyperlink{concepto-saas}{Software como Servicio}}, destinada a integrar información administrativa, operativa y de seguimiento dentro de un mismo entorno. La \textbf{\hyperlink{concepto-multitenant}{arquitectura multitenant}} y la \textbf{\hyperlink{concepto-separacion}{separación de información operativa}} permitirán orientar el diseño para que diferentes empresas puedan utilizar la plataforma manteniendo diferenciados sus registros y las funciones correspondientes a cada perfil.

De esta manera, la motivación principal del proyecto consiste en ofrecer una alternativa tecnológica que permita relacionar las operaciones de transporte, mantener información accesible para los participantes autorizados y apoyar la planificación de recorridos. TravelDLS no busca sustituir las decisiones administrativas de las empresas, sino proporcionar una herramienta que contribuya a organizar los procesos y facilitar la consulta de información durante la prestación del servicio.

La identificación del problema permite establecer una base para definir las características que deberá cumplir la solución. No obstante, las dificultades descritas a partir del análisis documental y los antecedentes deberán diferenciarse de aquellas que puedan comprobarse posteriormente mediante información obtenida directamente del entorno operativo.

En consecuencia, esta primera fase delimita las necesidades que motivan TravelDLS y proporciona los fundamentos para formular los objetivos de la solución, sin considerar como demostrados los beneficios que todavía requieren comprobación.


% =========================================================
% FASE 2
% =========================================================

\subsubsection{Fase 2. Definición de Objetivos de la Solución (Meta-requisitos)}
\label{sec:metodo-fase2}

La segunda fase de la metodología DSRM consiste en establecer las características que deberá cumplir la solución propuesta para responder a las necesidades identificadas durante la etapa anterior \parencite{peffers2007dsrm}. En TravelDLS, esta fase permite relacionar la problemática con los objetivos del proyecto y formular los meta-requisitos que orientarán el diseño de la plataforma.

Los meta-requisitos representan las capacidades generales que se espera que el sistema proporcione. Su formulación permite describir qué necesita resolver TravelDLS antes de establecer cómo serán implementadas técnicamente sus funcionalidades. Esta diferencia es importante porque permite separar las necesidades del sistema de las herramientas, estructuras y procedimientos que se utilizarán posteriormente para atenderlas.

A partir del problema identificado, uno de los principales meta-requisitos corresponde a la centralización de la información administrativa y operativa. La plataforma deberá permitir registrar y consultar los datos relacionados con empresas, clientes, conductores, vehículos y órdenes de transporte, manteniendo las asociaciones necesarias entre estos elementos.

Otro meta-requisito se relaciona con la administración de las solicitudes de transporte. TravelDLS deberá permitir que el cliente registre una solicitud, que la empresa establezca el precio correspondiente y que el cliente pueda aceptar o rechazar las condiciones del servicio. Asimismo, será necesario mantener un estado que permita identificar la situación actual de cada orden.

En relación con el seguimiento, la solución deberá permitir vincular los registros de ubicación y estado con las órdenes correspondientes, de manera que los usuarios autorizados puedan consultar información sobre el desplazamiento de sus cargas. Esta capacidad deberá considerar las limitaciones de conectividad y disponibilidad de los dispositivos utilizados.

La plataforma también deberá contemplar la asignación de recursos de transporte, permitiendo relacionar las operaciones con los conductores y vehículos correspondientes. Esta característica contribuirá a mantener organizada la información necesaria para identificar los recursos involucrados en cada servicio.

En cuanto a la planificación de rutas, TravelDLS deberá proporcionar un mecanismo que permita analizar recorridos con múltiples destinos. Para este componente se contempla la comparación de Vecino Más Cercano y 2-opt con Google OR-Tools, utilizando criterios equivalentes para estudiar sus resultados. La integración del módulo de optimización se orientará a apoyar la planificación, sin sustituir las decisiones que correspondan a la empresa transportista.

Asimismo, el sistema deberá mantener la separación de la información perteneciente a diferentes empresas transportistas, estableciendo las restricciones necesarias para que los usuarios accedan únicamente a los datos relacionados con sus funciones y servicios autorizados.

Finalmente, el proceso de contratación deberá mostrar el precio final del servicio antes de que el cliente lo acepte, incorporando el porcentaje destinado al mantenimiento de la plataforma. Esta condición responde al modelo de funcionamiento económico definido para TravelDLS y busca que el cliente conozca el monto correspondiente antes de confirmar la contratación.

Estos meta-requisitos constituyen una propuesta inicial que relaciona las necesidades identificadas con las capacidades esperadas del sistema. Su organización permitirá orientar las decisiones de diseño y establecer posteriormente los criterios necesarios para comprobar el funcionamiento de cada componente.

La siguiente matriz presenta esta relación inicial, sin asignar porcentajes de mejora o resultados cuantitativos que todavía no hayan sido determinados mediante procedimientos de evaluación.



% =========================================================
% MATRIZ DE META-REQUISITOS DE TRAVELDLS
% =========================================================

\begin{table}[H]
	\centering
	\caption{Matriz inicial de meta-requisitos de TravelDLS}
	\label{tab:meta-requisitos-traveldls}
	
	\small
	\setlength{\tabcolsep}{4.5pt}
	\renewcommand{\arraystretch}{1.45}
	
	\begin{tabularx}{\textwidth}{
			@{}
			>{\RaggedRight\arraybackslash}p{1.35cm}
			>{\RaggedRight\arraybackslash}p{3.05cm}
			>{\RaggedRight\arraybackslash}X
			>{\RaggedRight\arraybackslash}p{2.85cm}
			@{}
		}
		
		\toprule
		
		\textbf{Código} &
		\textbf{Necesidad identificada} &
		\textbf{Meta-requisito previsto} &
		\textbf{Relación con el diseño} \\
		
		\midrule
		
		MR-01 &
		Organización de información administrativa y operativa. &
		Registrar y consultar clientes, empresas, vehículos, conductores y órdenes desde una plataforma compartida. &
		Figuras~\ref{fig:diagrama-clases} y~\ref{fig:diagrama-er}. \\
		
		\midrule
		
		MR-02 &
		Administración de solicitudes de transporte. &
		Permitir la creación, cotización, aceptación, rechazo y consulta de estados de una orden. &
		Figuras~\ref{fig:diagrama-estados} y~\ref{fig:diagrama-secuencia}. \\
		
		\midrule
		
		MR-03 &
		Consulta de información sobre las cargas. &
		Asociar registros de ubicación y estado a las órdenes, permitiendo su consulta por usuarios autorizados. &
		Figuras~\ref{fig:diagrama-clases},~\ref{fig:diagrama-er} y~\ref{fig:diagrama-secuencia}. \\
		
		\midrule
		
		MR-04 &
		Coordinación de recursos de transporte. &
		Relacionar conductores y vehículos con las operaciones según las reglas establecidas. &
		Figuras~\ref{fig:diagrama-clases} y~\ref{fig:diagrama-er}. \\
		
		\midrule
		
		MR-05 &
		Planificación de recorridos con múltiples destinos. &
		Comparar métodos de optimización con entradas equivalentes y desarrollar el módulo previsto para el prototipo. &
		\textbf{\hyperref[sec:objetivos-especificos]{Objetivos del proyecto}} y módulo de optimización. \\
		
		\midrule
		
		MR-06 &
		Claridad en el precio del servicio. &
		Mostrar el monto final de la contratación, incluyendo el porcentaje de mantenimiento, antes de su aceptación. &
		Figuras~\ref{fig:diagrama-estados} y~\ref{fig:diagrama-secuencia}. \\
		
		\midrule
		
		MR-07 &
		Uso de la plataforma por diferentes empresas. &
		Mantener diferenciados los registros operativos de cada empresa y establecer controles de acceso según los perfiles. &
		Figuras~\ref{fig:diagrama-clases} y~\ref{fig:diagrama-er}. \\
		
		\bottomrule
		
	\end{tabularx}
	
	\vspace{0.2cm}
	
	\begin{minipage}{\textwidth}
		\footnotesize
		\textit{Nota.} Los meta-requisitos representan capacidades
		previstas para TravelDLS. Su cumplimiento deberá comprobarse
		posteriormente mediante criterios de aceptación y las
		evidencias correspondientes. La matriz no representa
		resultados de una evaluación finalizada.
	\end{minipage}
	
\end{table}

\subsubsection{Fase 3. Diseño y Desarrollo del Artefacto (Meta-diseños)}
\label{sec:metodo-fase3}

La tercera fase de la metodología DSRM comprende el diseño y la construcción de la solución tecnológica a partir de los objetivos establecidos previamente \parencite{peffers2007dsrm}. En TravelDLS, esta etapa permite relacionar los meta-requisitos con las decisiones de arquitectura, los modelos de información y los componentes que conformarán la plataforma web.

Los meta-diseños representan las decisiones técnicas destinadas a satisfacer las capacidades esperadas del sistema. En este proyecto, dichas decisiones se encuentran relacionadas con la organización de los módulos, el almacenamiento de información, las interacciones entre los usuarios y la definición de los procedimientos que permitirán desarrollar las operaciones de transporte.

Como parte del diseño actual de TravelDLS se dispone de un diagrama de clases y de un diagrama entidad-relación. El primero representa las principales entidades y servicios del sistema, mientras que el segundo permite comprender la estructura prevista para el almacenamiento de datos. Ambos modelos constituyen una base para establecer las relaciones necesarias entre usuarios, empresas, clientes, conductores, vehículos, órdenes, detalles de carga, pagos y registros de seguimiento.

El modelo de clases también permite reconocer los servicios destinados a administrar operaciones específicas, como la autenticación de usuarios, la gestión de órdenes y el procesamiento de pagos. Estas responsabilidades contribuyen a establecer una separación entre los componentes que representan los datos y aquellos que participan en la ejecución de determinadas funciones.

El diagrama de clases presentado en la Figura~\ref{fig:diagrama-clases} organiza los componentes principales del sistema alrededor de usuarios, roles, empresas, clientes, conductores, vehículos, órdenes, detalles de carga, pagos y seguimiento. Esta representación permite observar cómo los servicios de autenticación, órdenes y pagos se relacionan con las entidades del dominio, de manera que las operaciones previstas mantengan una separación entre responsabilidades funcionales y estructuras de información.

El diagrama entidad-relación de la Figura~\ref{fig:diagrama-er} complementa el modelo anterior al presentar las tablas, claves y asociaciones previstas para la persistencia de datos. Su incorporación permite relacionar el diseño conceptual con la estructura de almacenamiento, especialmente en los elementos asociados a empresas, clientes, conductores, vehículos, órdenes, detalles de carga, pagos y registros de ubicación.

En relación con el comportamiento del sistema, el diagrama de estados permite representar las situaciones que puede experimentar una orden de transporte. Su utilización contribuye a establecer las transiciones necesarias durante la creación, cotización, aceptación, ejecución y finalización de una solicitud.

Como se muestra en la Figura~\ref{fig:diagrama-estados}, la orden inicia en estado pendiente y avanza según las acciones realizadas por la empresa y el cliente. La asignación de precio, la aceptación o rechazo, la posible anulación, el pago, el inicio del viaje y la entrega establecen reglas que ayudan a mantener coherencia entre el avance administrativo de la solicitud y su ejecución operativa.

Por su parte, el diagrama de secuencia describe la interacción prevista entre el cliente, la empresa transportista, el conductor, el frontend, el backend y la base de datos. Este modelo permite observar cómo se intercambiará información durante los procesos de creación de órdenes, asignación de precios, registro de pagos, asignación de conductores y seguimiento de los servicios.

La Figura~\ref{fig:diagrama-secuencia} muestra el flujo principal de una orden desde su creación hasta el seguimiento del viaje. En este proceso, el cliente registra la solicitud, la empresa asigna el precio, el cliente acepta la cotización, se registra el pago, se asigna el conductor y posteriormente se consultan los registros de seguimiento. Esta secuencia permite vincular las operaciones de interfaz, backend y base de datos con los estados definidos para la orden.

La relación entre ambos diagramas es importante porque las operaciones representadas en las secuencias deben mantener coherencia con los estados definidos para las órdenes. Por esta razón, durante la consolidación del diseño será necesario revisar las reglas relacionadas con el momento de procesamiento del pago, la asignación del conductor y el inicio del traslado, evitando que un mismo proceso tenga interpretaciones diferentes entre los modelos.

En cuanto a la organización tecnológica, el diseño documentado contempla una plataforma web cuyo frontend utiliza Next.js y cuyo backend se desarrolla con AdonisJS. Estos componentes permiten separar la interfaz utilizada por los usuarios de los procedimientos encargados de procesar solicitudes, aplicar las reglas de negocio y administrar la información almacenada en la base de datos.

\paragraph{Diseño propuesto para la optimización de rutas}
El módulo de optimización de TravelDLS aún no está implementado. Como parte del meta-diseño se propone incorporar Google OR-Tools en una etapa posterior, con entradas y restricciones documentadas. Se distinguen dos ejercicios: un trayecto directo entre una recogida A y una entrega B, utilizado solo para ilustrar la elección de un camino en una red, y un recorrido que parte de A después de recoger la carga y visita varios destinos. Este segundo caso corresponde a la planificación con múltiples destinos prevista en los meta-requisitos y en el plan de ejecución.

Para el trayecto directo se propone representar la red vial mediante un grafo dirigido $G=(V,E)$, con un costo no negativo $c_{ij}$ para cada arco disponible. Una posible formulación en \texttt{SimpleMinCostFlow} asigna una unidad de suministro a A, una unidad de demanda a B, cero a los nodos intermedios y capacidad unitaria a cada arco \parencite{googleMinCostFlow}. La conservación del flujo se expresa como $\sum_j x_{vj}-\sum_i x_{iv}=b_v$, donde $b_A=1$, $b_B=-1$ y $b_v=0$ para los demás nodos; el objetivo es minimizar $\sum_{(i,j)\in E}c_{ij}x_{ij}$. Con costos enteros positivos, la unidad de flujo permite identificar un camino de costo mínimo. Las conexiones, sentidos de circulación y restricciones para vehículos de carga deberán obtenerse de datos viales adecuados; OR-Tools no los aporta por sí solo. Este ejercicio de camino mínimo no sustituye la comparación de métodos para múltiples destinos.

Para varias entregas se propone \texttt{RoutingModel}, con A como inicio fijo, una matriz común de distancias enteras y la obligación de visitar todos los destinos \parencite{googleTSP}. El ejercicio académico supondrá el retorno a A, conforme al TSP descrito en el marco teórico; antes de aplicar el modelo a un viaje real habrá que definir si el vehículo efectivamente regresa al origen y considerar las restricciones operativas pertinentes. Si se añaden recogidas y entregas emparejadas, deberán imponerse además las condiciones de mismo vehículo y precedencia, no solo registrar los pares mediante \texttt{AddPickupAndDelivery} \parencite{googlePickupDelivery}.

La comparación prevista para múltiples destinos utilizará la misma matriz, origen, destinos y condición de retorno para Vecino Más Cercano, su mejora mediante 2-opt y OR-Tools. Se registrarán el orden de visita, la distancia total y el tiempo de cálculo bajo condiciones de ejecución documentadas. Los ejercicios hipotéticos del apartado de resultados permiten verificar operaciones matemáticas, pero no constituyen pruebas de funcionamiento ni mediciones de rendimiento del módulo futuro. Tampoco se presupone que el solucionador de rutas encuentre siempre el óptimo global \parencite{googleTSP}.

\subsection{Metodología de Desarrollo}
\label{sec:metodologia-desarrollo}

La metodología de desarrollo de TravelDLS se integra en la fase 3 de la metodología de investigación en ciencia del diseño (Design Science Research Methodology, DSRM), denominada Diseño y Desarrollo del Artefacto. En esta fase, las decisiones de diseño y arquitectura derivadas de los meta-requisitos se utilizan como base para construir la plataforma web. Para organizar este proceso, se propone combinar Scrum y Extreme Programming (XP). Scrum se utilizará para planificar y distribuir el trabajo de forma incremental, mientras que XP orientará las prácticas de ingeniería hacia la calidad técnica del software \parencite{beck2000xp, schwaber2020scrum}.

Esta combinación se enmarca en el enfoque de investigación adoptado para TravelDLS. \textcite{peffers2007dsrm} describen la metodología DSRM como un proceso que comprende la identificación del problema, la definición de los objetivos de la solución, el diseño y desarrollo del artefacto, su demostración, su evaluación y la comunicación de los resultados. En consecuencia, el nivel táctico se concentra en la construcción del sistema a partir de los meta-requisitos establecidos, sin confundir la implementación de las funcionalidades con los resultados que deberán obtenerse durante su evaluación.

La metodología de desarrollo describe cómo se organizará la construcción del artefacto TravelDLS a partir de los siete meta-requisitos definidos en la Tabla~\ref{tab:meta-requisitos-traveldls}. Esta actividad se relaciona con la fase de diseño y desarrollo del artefacto dentro de la metodología de investigación en ciencia del diseño (DSRM).\textcite{peffers2007dsrm} presentan la DSRM como un proceso que vincula la identificación del problema, los objetivos de la solución, el diseño y desarrollo, la demostración, la evaluación y la comunicación de resultados. Por ello, las matrices de esta sección conectan las necesidades del sistema con las decisiones de diseño, las historias de usuario y las formas previstas de comprobar su cumplimiento.

Para organizar el trabajo se propone Scrum como marco de gestión y se consideran prácticas técnicas de Extreme Programming (XP) para apoyar la calidad del código. \textcite{bahit2012scrum} explica cómo pueden combinarse ambos enfoques, utilizando Scrum para organizar el trabajo en iteraciones y XP para reforzar las prácticas de desarrollo. La aplicación de estas prácticas en TravelDLS deberá verificarse durante la implementación y mediante las pruebas correspondientes \parencite{beck2000xp, bahit2012scrum, schwaber2020scrum}.

\subsubsection{Gestión del desarrollo mediante Scrum}

Scrum es un marco de trabajo que permite abordar problemas complejos mediante un proceso iterativo e incremental. Se fundamenta en el empirismo y el pensamiento Lean, y establece la transparencia, la inspección y la adaptación como pilares para gestionar el trabajo. Estos principios permiten revisar el avance, identificar desviaciones y ajustar las actividades según la información obtenida durante el desarrollo \parencite{schwaber2020scrum}.

Para TravelDLS, se propone utilizar Scrum para organizar las actividades de desarrollo mediante un Product Backlog, historias de usuario y Sprints. El Product Backlog reunirá las funcionalidades pendientes del sistema y se actualizará conforme se definan nuevas necesidades o se precise el alcance de las ya identificadas. Las historias de usuario describirán las funcionalidades desde la perspectiva de los perfiles que utilizarán la plataforma, manteniendo su relación con los meta-requisitos establecidos. El Product Backlog reúne las historias de usuario priorizadas según su valor funcional y las dependencias entre tareas. En cada sprint, el equipo selecciona el trabajo que abordará y revisa el incremento obtenido; la retrospectiva permite identificar ajustes para la siguiente iteración (Bahit, 2012; Schwaber \& Sutherland, 2020). La asignación de historias a los sprints es provisional y debe concordar con el \mbox{cronograma vigente}.

El trabajo se organizará en Sprints de dos a tres semanas, según la planificación propuesta en el documento metodológico. Al inicio de cada Sprint se establecerá un objetivo y se seleccionarán los elementos del Product Backlog que se atenderán, considerando su prioridad, complejidad y dependencias técnicas. Durante la ejecución se desarrollarán las funcionalidades previstas y se realizarán las comprobaciones correspondientes. Al finalizar cada Sprint, se revisará el incremento obtenido y se identificarán ajustes para las iteraciones siguientes, mediante actividades de revisión y retrospectiva (Schwaber \& Sutherland, 2020).

La planificación tendrá en cuenta la arquitectura tecnológica definida para TravelDLS: Next.js para el frontend, AdonisJS para el backend y PostgreSQL para la persistencia de datos. El backend se organizará mediante el patrón MVC y una capa de servicios para separar las responsabilidades de los controladores y la lógica de negocio. La distribución de actividades considerará las dependencias entre los componentes, de modo que la implementación de las interfaces se integre con los servicios y las estructuras de datos correspondientes.

\subsubsection{Prácticas de ingeniería mediante Extreme Programming (XP)}

Extreme Programming (XP) propone prácticas de ingeniería que buscan mejorar la calidad del software y facilitar su adaptación durante el desarrollo. \textcite{beck2000xp} describe este enfoque a partir de prácticas que promueven la retroalimentación frecuente, la simplicidad del diseño y la colaboración entre los integrantes del equipo. En TravelDLS se propone incorporar cuatro prácticas: desarrollo guiado por pruebas, refactorización continua, programación en pares y propiedad colectiva del código, e integración continua.

\begin{itemize}
	\item \textbf{Desarrollo guiado por pruebas (TDD):} El desarrollo guiado por pruebas (\textit{Test-Driven Development}, TDD) consiste en escribir una prueba automatizada antes del código funcional que permitirá satisfacerla. A partir del resultado de la prueba, se implementa la funcionalidad y se realizan los ajustes necesarios hasta que se cumpla el comportamiento esperado \parencite{beck2000xp}. En TravelDLS se propone aplicar esta práctica principalmente a las reglas de negocio del backend, como las validaciones de usuarios, la asignación de vehículos y conductores, la gestión de órdenes y el procesamiento de información relacionada con los recorridos. Las pruebas permitirán comprobar el comportamiento esperado de cada componente y detectar errores durante su implementación.
\end{itemize}

\begin{itemize}
	\item \textbf{Refactorización continua:} La refactorización consiste en mejorar la estructura interna del código sin modificar su comportamiento observable \parencite{beck2000xp}. En TravelDLS se propone revisar y reorganizar el código conforme se incorporen nuevas funcionalidades, especialmente cuando existan instrucciones repetidas, responsabilidades mezcladas o procedimientos que puedan simplificarse. Esta práctica se aplicará a los componentes del frontend desarrollado con Next.js y a los módulos del backend construido con AdonisJS, procurando conservar la separación entre la interfaz, los controladores, la capa de servicios y el acceso a los datos.
	
	\item \textbf{Programación en pares y propiedad colectiva del código:} La programación en pares consiste en que dos integrantes colaboren en una misma tarea de programación, mientras que la propiedad colectiva permite que los integrantes del equipo puedan revisar y modificar diferentes partes del código, en lugar de limitar cada componente a una sola persona \parencite{beck2000xp}. Para TravelDLS se propone utilizar estas prácticas en funcionalidades que involucren varios componentes del sistema, como la autenticación, la gestión de órdenes y la optimización de recorridos. La colaboración facilitará el intercambio de conocimientos, la revisión de las soluciones implementadas y el mantenimiento del código. Su aplicación se ajustará a la disponibilidad de los integrantes y a la distribución de las tareas de cada sprint.
\end{itemize}

\begin{itemize}
	\item \textbf{Integración continua:} \textcite{bahit2012scrum} explica que el objetivo de la integración continua es incorporar las nuevas funcionalidades a las existentes, procurando que los cambios realizados no afecten el funcionamiento de las funcionalidades anteriores. En TravelDLS, se propone aplicar esta práctica mediante la integración progresiva de los cambios en el código y la ejecución de pruebas que permitan comprobar que las nuevas funcionalidades se incorporen correctamente sin alterar las ya desarrolladas. Esto contribuirá a detectar errores durante el desarrollo y a mantener el funcionamiento conjunto de los componentes del sistema.
\end{itemize}

\subsubsection{Matriz de meta-diseños}

Los meta-diseños resumen las decisiones técnicas previstas para responder a los meta-requisitos; su cumplimiento se comprobará durante el desarrollo.


\begingroup
	\small
	\setlength{\tabcolsep}{4.5pt}
	\renewcommand{\arraystretch}{1.45}
	\begin{longtable}{@{}>{\RaggedRight\arraybackslash}p{1.5cm}>{\RaggedRight\arraybackslash}p{1.8cm}>{\RaggedRight\arraybackslash}p{4.1cm}>{\RaggedRight\arraybackslash}p{2.3cm}>{\RaggedRight\arraybackslash}p{5.3cm}@{}}
		\caption{Matriz de meta-diseños asociados a los meta-requisitos de TravelDLS}\label{tab:meta_disenos}\\
		\toprule
		\textbf{Código} & \textbf{MR asociado} & \textbf{Decisión de diseño} & \textbf{Historias} & \textbf{Verificación} \\
		\midrule
		\endfirsthead
		\caption[]{Matriz de meta-diseños asociados a los meta-requisitos de TravelDLS (continuación)}\\
		\toprule
		\textbf{Código} & \textbf{MR asociado} & \textbf{Decisión de diseño} & \textbf{Historias} & \textbf{Verificación} \\
		\midrule
		\endhead
		\endfoot
		\bottomrule
		\endlastfoot
		MD-01 & MR-01 & Modelo relacional en PostgreSQL para usuarios, roles, empresas, clientes, conductores, vehículos, órdenes y registros asociados; operaciones CRUD mediante AdonisJS, organizadas con MVC y capa de servicios. & HU-01, HU-05, HU-14, HU-29 & Integridad de relaciones; persistencia de altas y cambios; consulta paginada; errores de validación. \\ \midrule
		MD-02 & MR-02 & Flujo de órdenes con estados definidos y operaciones para crear solicitud, cotizar, aceptar o rechazar, consultar y modificar según el estado permitido. & HU-16, HU-20 & Transiciones válidas; cotización visible antes de decidir; rechazo de operaciones incompatibles. \\ \midrule
		MD-03 & MR-03 & Registros de ubicación y estado vinculados a una orden; consulta restringida a participantes autorizados; mapa y ubicación real separados de la simulación del recorrido. & HU-22, HU-24 & Última ubicación asociada al viaje correcto; permisos aplicados; simulación sin modificar GPS ni estado real. \\ \midrule
		MD-04 & MR-04 & Relaciones entre empresas, conductores, camiones y pedidos, con reglas de disponibilidad, asignación, estado operativo y reporte de mantenimiento. & HU-11, HU-15, HU-21, HU-22 & No permitir asignaciones incompatibles; conductor consulta solo su vehículo/viajes; cambios de disponibilidad coherentes. \\ \midrule
		MD-05 & MR-05 & Módulo de recorridos con origen y destinos, matriz de costo, integración de Google OR-Tools y comparación con Vecino Más Cercano y 2-opt; visualización mediante Mapbox. & HU-30, HU-33 & Mismas entradas y condiciones en la comparación; costo y tiempo de cálculo registrados; no afirmar óptimo garantizado. \\ \midrule
		MD-06 & MR-06 & Presentación del precio final antes de la aceptación, incluyendo el porcentaje de mantenimiento previsto en el modelo económico; flujo de pago del prototipo identificado explícitamente como simulado. & HU-16, HU-19, HU-20, HU-25, HU-28 & Importe final visible antes de aceptar; reglas de estado y total consistentes; ninguna operación se presenta como cobro bancario real. \\ \midrule
		MD-07 & MR-07 & Autenticación, sesión y autorización por rol, con comprobación de pertenencia de los registros a la empresa, cliente, conductor o pedido correspondiente. & HU-03, HU-04, HU-06, HU-07, HU-11, HU-14, HU-16, HU-18, HU-21, HU-23, HU-26, HU-29, HU-30, HU-33 & Denegar acceso directo o por API a datos ajenos; mantener separada la información por empresa y perfil. \\ \midrule
	\end{longtable}
\endgroup

\textit{Nota.} Elaboración propia a partir de los meta-requisitos, los diagramas de diseño y las historias de usuario de TravelDLS.

\subsubsection{Matriz de trazabilidad de historias de usuario}

La Tabla~\ref{tab:trazabilidad_tactica} muestra la contribución de cada historia de usuario a los meta-requisitos de la Tabla~\ref{tab:meta-requisitos-traveldls}.


\begingroup
	\small
	\setlength{\tabcolsep}{4.5pt}
	\renewcommand{\arraystretch}{1.45}
	\begin{longtable}{@{}>{\RaggedRight\arraybackslash}p{1.7cm}>{\RaggedRight\arraybackslash}p{3.9cm}>{\RaggedRight\arraybackslash}p{2.8cm}>{\RaggedRight\arraybackslash}p{7cm}@{}}
		\caption{Matriz de trazabilidad entre historias de usuario y meta-requisitos}\label{tab:trazabilidad_tactica}\\
		\toprule
		\textbf{Historia} & \textbf{Resumen} & \textbf{Meta-requisitos} & \textbf{Área funcional} \\
		\midrule
		\endfirsthead
		\caption[]{Matriz de trazabilidad entre historias de usuario y meta-requisitos (continuación)}\\
		\toprule
		\textbf{Historia} & \textbf{Resumen} & \textbf{Meta-requisitos} & \textbf{Área funcional} \\
		\midrule
		\endhead
		\endfoot
		\bottomrule
		\endlastfoot
		HU-01 & Cimientos de base de datos y backend & MR-01 & Preparación de infraestructura y persistencia \\ \midrule
		HU-02 & Registrarse como cliente & MR-01, MR-07 & Registro de clientes y asignación del rol cliente \\ \midrule
		HU-03 & Iniciar, conservar y cerrar sesión & MR-01, MR-07 & Autenticación y ciclo de sesión \\ \midrule
		HU-04 & Acceder a navegación y operaciones por rol & MR-07 & Navegación y autorización por perfil \\ \midrule
		HU-05 & Registrar usuarios con rol desde administración & MR-01, MR-07 & Administración de usuarios \\ \midrule
		HU-06 & Administrar empresas & MR-01, MR-07 & Gestión de empresas \\ \midrule
		HU-07 & Consultar y dar de baja clientes desde administración & MR-01, MR-07 & Consulta y baja de clientes \\ \midrule
		HU-08 & Completar y actualizar perfil de cliente & MR-01, MR-07 & Perfil de cliente \\ \midrule
		HU-09 & Actualizar perfil y tarifa de empresa & MR-01, MR-06, MR-07 & Perfil y tarifa de empresa \\ \midrule
		HU-10 & Completar y actualizar perfil de conductor & MR-01, MR-07 & Perfil de conductor \\ \midrule
		HU-11 & Registrar, vincular y administrar conductores & MR-01, MR-04, MR-07 & Gestión y vinculación de conductores \\ \midrule
		HU-12 & Administrar categorías de vehículos & MR-01, MR-04 & Categorías de vehículos \\ \midrule
		HU-13 & Administrar camiones y asignarlos a conductores & MR-01, MR-04, MR-07 & Gestión de camiones y asignación de conductores \\ \midrule
		HU-14 & Consultar camión asignado & MR-01, MR-04, MR-07 & Consulta del camión asignado \\ \midrule
		HU-15 & Reportar mantenimiento del camión & MR-04 & Reporte de mantenimiento \\ \midrule
		HU-16 & Solicitar transporte con empresa, carga y direcciones & MR-02, MR-07 & Solicitud de transporte \\ \midrule
		HU-17 & Consultar y gestionar pedidos del cliente & MR-02, MR-07 & Consulta y cambios permitidos de pedidos \\ \midrule
		HU-18 & Consultar solicitudes recibidas por empresa & MR-02, MR-07 & Consulta de solicitudes por empresa \\ \midrule
		HU-19 & Establecer el precio de un pedido & MR-02, MR-07 & Cotización del servicio \\ \midrule
		HU-20 & Responder a la cotización de transporte & MR-02, MR-07 & Aceptación o rechazo de cotización \\ \midrule
		HU-21 & Asignar conductor a un pedido & MR-04, MR-07 & Asignación de conductor al pedido \\ \midrule
		HU-22 & Consultar, iniciar y finalizar viajes & MR-03, MR-04, MR-07 & Gestión de viajes \\ \midrule
		HU-23 & Compartir y consultar ubicación del viaje & MR-03, MR-07 & Ubicación real y consulta del seguimiento \\ \midrule
		HU-24 & Visualizar simulación del recorrido & MR-03, MR-07 & Simulación visual separada del seguimiento real \\ \midrule
		HU-25 & Realizar pago simulado por tarjeta o transferencia & MR-02, MR-07 & Pago simulado de una cotización \\ \midrule
		HU-26 & Consultar historial y detalle de pagos como cliente & MR-01, MR-07 & Historial y detalle de pagos del cliente \\ \midrule
		HU-27 & Consultar pagos recibidos como empresa & MR-01, MR-07 & Consulta de pagos por empresa \\ \midrule
		HU-28 & Cancelar pago simulado con motivo & MR-02, MR-07 & Cancelación permitida de pago simulado \\ \midrule
		HU-29 & Consultar indicadores de los cuatro perfiles & MR-01, MR-07 & Indicadores filtrados por perfil \\ \midrule
		HU-30 & Definir un recorrido con múltiples destinos & MR-05, MR-07 & Definición de recorrido con múltiples destinos \\ \midrule
		HU-31 & Calcular una ruta optimizada & MR-05, MR-07 & Cálculo de una ruta propuesta \\ \midrule
		HU-32 & Visualizar y guardar la ruta propuesta & MR-03, MR-07 & Visualización y persistencia de ruta planificada \\ \midrule
		HU-33 & Comparar los métodos de optimización & MR-05, MR-07 & Comparación experimental de métodos \\ \midrule
	\end{longtable}
\endgroup
\textit{Nota.} Elaboración propia a partir de las historias de usuario de TravelDLS. Los vínculos MR–HU son una propuesta de trazabilidad elaborada a partir de los resúmenes y descripciones disponibles.

\subsubsection{Matriz de criterios de aceptación}

Los criterios siguientes sintetizan las condiciones previstas para las historias de usuario de TravelDLS. Constituyen condiciones previstas para comprobar el comportamiento de las historias y no resultados de pruebas ya ejecutadas. En la implementación, los criterios deberán probarse con casos válidos, datos inválidos, permisos insuficientes y escenarios de error cuando corresponda.


\begingroup
	\small
	\setlength{\tabcolsep}{4.5pt}
	\renewcommand{\arraystretch}{1.45}
	\begin{longtable}{@{}>{\RaggedRight\arraybackslash}p{1.7cm}>{\RaggedRight\arraybackslash}p{3.9cm}>{\RaggedRight\arraybackslash}p{10.1cm}@{}}
		\caption{Matriz de criterios de aceptación}\label{tab:criterios_tactica}\\
		\toprule
		\textbf{Historia} & \textbf{Funcionalidad} & \textbf{Criterios de aceptación propuestos} \\
		\midrule
		\endfirsthead
		\caption[]{Matriz de criterios de aceptación (continuación)}\\
		\toprule
		\textbf{Historia} & \textbf{Funcionalidad} & \textbf{Criterios de aceptación propuestos} \\
		\midrule
		\endhead
		\endfoot
		\bottomrule
		\endlastfoot
		HU-01 & Cimientos de base de datos y backend & Preparar el backend y la persistencia de la aplicación. Configurar AdonisJS y PostgreSQL. Crear migraciones y modelos con sus relaciones. Verificar arranque, conexión y operaciones básicas. \\ \midrule
		HU-02 & Registrarse como cliente & Crear una cuenta con nombre, correo y contraseña válidos, siempre con rol cliente. Informar campos inválidos o correo duplicado sin crear cuentas repetidas. Acceder al área de cliente al completar el registro. \\ \midrule
		HU-03 & Iniciar, conservar y cerrar sesión & Acceder con credenciales válidas y mostrar error si son incorrectas. Conservar la sesión al recargar y solicitar acceso cuando expire. Cerrar sesión e impedir volver a consultar contenido privado. \\ \midrule
		HU-04 & Acceder a navegación y operaciones por rol & Mostrar las áreas y acciones correspondientes a cada rol. Impedir el acceso directo a áreas de otro perfil. Rechazar operaciones no autorizadas también desde la API. \\ \midrule
		HU-05 & Registrar usuarios con rol desde administración & Registrar nombre, correo, contraseña y un rol permitido. Informar datos inválidos o correo duplicado. Guardar el usuario con el rol elegido y confirmar el registro. \\ \midrule
		HU-06 & Administrar empresas & Consultar y buscar empresas en un listado paginado. Registrar y actualizar sus datos, usuario asociado y fotografía. Confirmar la baja antes de retirar una empresa del listado. \\ \midrule
		HU-07 & Consultar y dar de baja clientes desde administración & Consultar clientes y sus datos en un listado paginado. Solicitar confirmación antes de dar de baja a un cliente. Actualizar el listado tras la baja y mostrar errores si falla. \\ \midrule
		HU-08 & Completar y actualizar perfil de cliente & Consultar y completar los datos del perfil propio. Actualizar datos y fotografía, validando la información ingresada. Conservar los cambios después de recargar. \\ \midrule
		HU-09 & Actualizar perfil y tarifa de empresa & Consultar los datos de la empresa asociada a la sesión. Actualizar datos, fotografía y tarifa por kilogramo. Conservar los cambios e informar valores inválidos. \\ \midrule
		HU-10 & Completar y actualizar perfil de conductor & Consultar y completar el perfil del conductor. Actualizar datos y fotografía respetando las validaciones de identificación y licencia. Conservar los cambios y mostrar errores de validación. \\ \midrule
		HU-11 & Registrar, vincular y administrar conductores & Registrar conductores y vincularlos a una empresa según los permisos del usuario. Consultar y mantener los conductores mediante las acciones habilitadas para cada rol. Rechazar datos inválidos y operaciones sobre empresas no autorizadas. \\ \midrule
		HU-12 & Administrar categorías de vehículos & Consultar categorías y su estado en un listado paginado. Crear y editar categorías con datos válidos. Confirmar la baja y actualizar las opciones de los formularios de camiones. \\ \midrule
		HU-13 & Administrar camiones y asignarlos a conductores & Registrar y editar placa, chasis, categoría, empresa, conductor y estado. Impedir asignar un conductor que ya tenga otro camión activo. Consultar la flota y confirmar la baja de un camión. \\ \midrule
		HU-14 & Consultar camión asignado & Mostrar los datos del camión asignado al conductor autenticado. Indicar cuando no existe un camión asignado. Mostrar los errores de consulta sin presentar datos de otro conductor. \\ \midrule
		HU-15 & Reportar mantenimiento del camión & Reportar mantenimiento únicamente del camión asignado. Guardar el estado de mantenimiento y mostrarlo al recargar. Informar si no existe asignación o si falla la operación. \\ \midrule
		HU-16 & Solicitar transporte con empresa, carga y direcciones & Seleccionar empresa e indicar carga, cantidad, peso, embalaje y direcciones. Guardar el pedido y sus detalles, validando los datos obligatorios. Informar fallos o registros parciales sin duplicar la solicitud. \\ \midrule
		HU-17 & Consultar y gestionar pedidos del cliente & Consultar pedidos propios con estado, importe y detalle. Permitir cambios y bajas únicamente en los estados habilitados. Actualizar los datos después de cada operación y respetar filtros y paginación. \\ \midrule
		HU-18 & Consultar solicitudes recibidas por empresa & Consultar las solicitudes correspondientes a la empresa. Mostrar cliente, carga, estado e importe con filtros y paginación. Reflejar cambios de estado e informar cuando no existan solicitudes. \\ \midrule
		HU-19 & Establecer el precio de un pedido & Asignar un importe válido a un pedido de la empresa. Dejar el pedido pendiente de aprobación del cliente. Mostrar al cliente el precio actualizado e informar si la operación falla. \\ \midrule
		HU-20 & Responder a la cotización de transporte & Consultar el precio propuesto antes de decidir. Rechazar la cotización y reflejar el estado denegado. Confirmar la aceptación mediante el flujo de pago simulado, sin cambios de estado incompatibles. \\ \midrule
		HU-21 & Asignar conductor a un pedido & Seleccionar un conductor autorizado para el pedido. Rechazar la asignación al alcanzar cinco viajes activos según la regla vigente. Mostrar el pedido entre los viajes del conductor asignado. \\ \midrule
		HU-22 & Consultar, iniciar y finalizar viajes & Consultar únicamente los viajes asignados al conductor. Registrar inicio y entrega, actualizando la disponibilidad del conductor. Impedir iniciar pedidos anulados y detener el seguimiento al finalizar. \\ \midrule
		HU-23 & Compartir y consultar ubicación del viaje & Compartir la ubicación del conductor durante el viaje. Permitir que cliente y empresa participantes consulten la última posición. Informar falta de señal, permiso GPS o ubicación, sin exponer viajes ajenos. \\ \midrule
		HU-24 & Visualizar simulación del recorrido & Identificar el recorrido como una simulación, separada de la ubicación real. Permitir iniciar, pausar y reiniciar la animación con progreso y tiempo estimado. No modificar la ubicación real ni el estado del pedido. \\ \midrule
		HU-25 & Realizar pago simulado por tarjeta o transferencia & Simular el pago de un pedido propio pendiente de aprobación y con importe válido. Validar tarjeta o transferencia e impedir pagos duplicados. Guardar la referencia y aceptar el pedido, indicando que no existe un cobro real. \\ \midrule
		HU-26 & Consultar historial y detalle de pagos como cliente & Consultar pagos propios con estado, método, importe y referencia. Filtrar, paginar y abrir el detalle de un pago. Mostrar totales según el alcance consultado y gestionar listas vacías o errores. \\ \midrule
		HU-27 & Consultar pagos recibidos como empresa & Consultar únicamente los pagos asociados a la empresa. Filtrar, paginar y abrir el detalle de las operaciones. Mostrar importes y estados de los pagos simulados. \\ \midrule
		HU-28 & Cancelar pago simulado con motivo & Solicitar motivo válido y confirmación para cancelar un pago propio. Impedir cancelar pagos ya cancelados o de pedidos en tránsito o entregados. Guardar la cancelación y anular el pedido si estaba aceptado. \\ \midrule
		HU-29 & Consultar indicadores de los cuatro perfiles & Mostrar a cada perfil sus indicadores y accesos rápidos. Calcular los valores con datos reales y respetar el alcance de la consulta. Distinguir ausencia de actividad de errores de consulta y permitir vista móvil. \\ \midrule
		HU-30 & Definir un recorrido con múltiples destinos & Permitir agregar, editar y eliminar destinos antes de calcular la ruta. Validar coordenadas y datos obligatorios, indicando puntos inválidos o repetidos. Guardar el recorrido asociado a la empresa y limitar su acceso a usuarios autorizados. \\ \midrule
		HU-31 & Calcular una ruta optimizada & Visitar todos los destinos válidos respetando el origen y la condición de finalización definidos. Mostrar el orden recomendado y el costo total, con la unidad utilizada. Informar entradas incompletas, destinos inaccesibles o ausencia de solución sin mostrar resultados engañosos. \\ \midrule
		HU-32 & Visualizar y guardar la ruta propuesta & Representar el recorrido y los destinos numerados en el orden calculado. Guardar la propuesta únicamente tras la confirmación del usuario y permitir consultarla nuevamente. Distinguir la ruta planificada de la ubicación real y de la simulación del recorrido. \\ \midrule
		HU-33 & Comparar los métodos de optimización & Comparar Vecino Más Cercano, su mejora con 2-opt y OR-Tools con los mismos puntos, matriz y condiciones. Mostrar por método el recorrido, costo total y tiempo de cálculo. Conservar entradas, parámetros y resultados para repetir la evaluación, sin asumir un método ganador. \\ \midrule
	\end{longtable}
\endgroup
\textit{Nota.} HU-01 reúne las tareas iniciales de configuración y comprobación de la persistencia.

\subsubsection{Planificación propuesta de los sprints}

La Tabla \ref{tab:sprints} presenta una agrupación funcional inicial de las historias de usuario. Las fechas y la duración de los sprints deben validarse con el cronograma general antes de establecerse como planificación definitiva.


\begingroup
	\small
	\setlength{\tabcolsep}{4.5pt}
	\renewcommand{\arraystretch}{1.45}
	\begin{longtable}{@{}>{\RaggedRight\arraybackslash}p{1.9cm}>{\RaggedRight\arraybackslash}p{3.45cm}>{\RaggedRight\arraybackslash}p{2.6cm}>{\RaggedRight\arraybackslash}p{7.45cm}@{}}
		\caption{Planificación propuesta de los sprints}\label{tab:sprints}\\
		\toprule
		\textbf{Iteración} & \textbf{Objetivo} & \textbf{Historias} & \textbf{Entregable previsto} \\
		\midrule
		\endfirsthead
		\caption[]{Planificación propuesta de los sprints (continuación)}\\
		\toprule
		\textbf{Iteración} & \textbf{Objetivo} & \textbf{Historias} & \textbf{Entregable previsto} \\
		\midrule
		\endhead
		\endfoot
		\bottomrule
		\endlastfoot
		Sprint 1 & Fundamentos y acceso & HU-01 a HU-04 & Backend y PostgreSQL configurados; registro, autenticación, sesión y control de acceso básico comprobados. \\ \midrule
		Sprint 2 & Administración y perfiles & HU-05 a HU-10 & Gestión de usuarios, empresas y clientes; actualización de perfiles con validaciones. \\ \midrule
		Sprint 3 & Conductores y flota & HU-11 a HU-15 & Gestión de conductores, categorías y camiones; asignaciones, disponibilidad y reporte de mantenimiento. \\ \midrule
		Sprint 4 & Pedidos y cotizaciones & HU-16 a HU-20 & Solicitud, consulta, cotización y respuesta del cliente; reglas de estados y visualización del importe. \\ \midrule
		Sprint 5 & Asignación, viajes y seguimiento & HU-21 a HU-24 & Asignación de conductor, inicio y finalización de viajes, ubicación real y simulación separada. \\ \midrule
		Sprint 6 & Pagos simulados e indicadores & HU-25 a HU-29 & Pago simulado, historial y consultas por perfil, cancelación permitida e indicadores con permisos. \\ \midrule
		Sprint 7 & Planificación y comparación de rutas & HU-30 a HU-33 & Definición de destinos, cálculo, mapa, persistencia y comparación de métodos con entradas equivalentes. \\ \midrule
	\end{longtable}
\endgroup
\textit{Nota.} Elaboración propia.


\subsubsection{Construcción del prototipo}

La construcción de TravelDLS se organiza mediante las historias de usuario y los sprints propuestos. En el backend se desarrollaron rutas y servicios para autenticación, administración de usuarios y recursos de transporte, gestión de órdenes, cotización, asignación de conductores, pagos simulados, consulta de ubicación y confirmación de entregas. En el frontend se incorporaron interfaces para administración, clientes, empresas transportistas y conductores, con pantallas para consultar recursos y pedidos, solicitar transporte, responder a cotizaciones y consultar información del viaje. La disponibilidad de estos componentes no supone que todos los recorridos entre interfaz y servicios se hayan verificado.

Para el seguimiento se incorporó el envío de la ubicación reportada por el dispositivo del conductor y su consulta por participantes autorizados. La visualización utiliza OpenStreetMap y Leaflet, y durante el desarrollo el trazado de rutas puede consultar un servicio de prueba basado en OSRM. El backend conserva la última posición asociada al pedido, pero aún no construye un historial de puntos. La incorporación de Mapbox y del historial requiere adaptar las vistas y definir el almacenamiento de ubicaciones sucesivas. También se incorporó código para gestionar con Cloudinary fotografías de perfiles, empresas y entregas; las pruebas unitarias validan los archivos antes de su carga, pero falta comprobar el envío al servicio externo.

En la contratación se desarrollaron el registro de precios de las cargas, la respuesta del cliente a la cotización y un pago simulado por tarjeta o transferencia. Las pruebas unitarias comprueban el cálculo de montos y algunas validaciones del medio de pago; una prueba funcional comprueba el registro de una transferencia simulada y el rechazo de un segundo pago para la misma orden. Sigue pendiente verificar el recargo del 8\,\% y la secuencia completa de cotización, aceptación y pago. Se dispone de documentación técnica del frontend, los mapas y la confirmación de entregas, así como de archivos de prueba para contratos del frontend; los manuales de usuario todavía no se han documentado.

\subsection{Cronograma de actividades}

El cronograma de trabajo presenta las actividades desarrolladas durante
el proceso de documentación y desarrollo del software TravelDLS,
indicando el período establecido para cada una y su estado de avance.

\begin{center}
	\begin{ganttchart}[
		hgrid,
		vgrid,
		x unit=1.35cm,
		y unit title=0.60cm,
		y unit chart=0.42cm,
		title height=1,
		bar height=0.24,
		group height=0.28,
		title label font=\scriptsize,
		bar label font=\scriptsize,
		group label font=\bfseries\scriptsize,
		title/.style={
			fill=white,
			draw=grislinea
		},
		group/.style={
			fill=white,
			draw=none
		},
		bar/.style={
			draw=none
		},
		bar label node/.append style={
			align=right,
			text width=6.1cm,
			anchor=east,
			xshift=-0.2cm
		},
		group label node/.append style={
			align=right,
			text width=6.1cm,
			anchor=east,
			xshift=-0.2cm
		}
		]{1}{7}
		
		% ENCABEZADO DE SEMANAS
		\gantttitle{\shortstack{\scriptsize Semana 1\\\tiny 7--11 sep}}{1}
		\gantttitle{\shortstack{\scriptsize Semana 2\\\tiny 14--18 sep}}{1}
		\gantttitle{\shortstack{\scriptsize Semana 3\\\tiny 21--25 sep}}{1}
		\gantttitle{\shortstack{\scriptsize Semana 4\\\tiny 28 sep--2 oct}}{1}
		\gantttitle{\shortstack{\scriptsize Semana 5\\\tiny 5--9 oct}}{1}
		\gantttitle{\shortstack{\scriptsize Semana 6\\\tiny 12--16 oct}}{1}
		\gantttitle{\shortstack{\scriptsize Semana 7\\\tiny 19--23 oct}}{1}\\
		
		% DOCUMENTACIÓN
		\ganttgroup{DOCUMENTACIÓN}{1}{7}\\
		
		\ganttbar[bar/.style={fill=completado,draw=none}]
		{Revisión bibliográfica}{1}{1}\\
		
		\ganttbar[bar/.style={fill=completado,draw=none}]
		{Objetivos y tema}{1}{1}\\
		
		\ganttbar[bar/.style={fill=completado,draw=none}]
		{Planteamiento, alcance y limitaciones}{2}{2}\\
		
		\ganttbar[bar/.style={fill=completado,draw=none}]
		{Marco teórico}{2}{2}\\
		
		\ganttbar[bar/.style={fill=completado,draw=none}]
		{Factibilidad del proyecto}{3}{3}\\
		
		\ganttbar[bar/.style={fill=completado,draw=none}]
		{Metodología}{4}{6}\\
		
		\ganttbar[bar/.style={fill=proceso,draw=none}]
		{Plan de ejecución}{6}{6}\\
		
		\ganttbar[bar/.style={fill=proceso,draw=none}]
		{Resultados del plan de ejecución}{7}{7}\\
		
		% BACKEND
		\ganttgroup{BACKEND}{1}{6}\\
		
		\ganttbar[bar/.style={fill=completado,draw=none}]
		{Diagrama entidad-relación (DER)}{1}{1}\\
		
		\ganttbar[bar/.style={fill=completado,draw=none}]
		{Usuarios, roles y compañía -- CRUD}{2}{2}\\
		
		\ganttbar[bar/.style={fill=completado,draw=none}]
		{Clientes, conductores y camiones -- CRUD}{3}{3}\\
		
		\ganttbar[bar/.style={fill=completado,draw=none}]
		{Categorías de camiones y órdenes -- CRUD}{4}{4}\\
		
		\ganttbar[bar/.style={fill=completado,draw=none}]
		{Autenticación JWT y pruebas de integración}{4}{4}\\
		
		\ganttbar[bar/.style={fill=completado,draw=none}]
		{Migraciones, modelos y validaciones}{5}{5}\\
		
		\ganttbar[bar/.style={fill=proceso,draw=none}]
		{Finalización de CRUD y pruebas del backend}{6}{6}\\
		
		% FRONTEND
		\ganttgroup{FRONTEND}{7}{7}\\
		
		\ganttbar[bar/.style={fill=proceso,draw=none}]
		{Configuración de Next.js y autenticación JWT}{7}{7}
		
	\end{ganttchart}
	
	\vspace{0.3cm}
	
	\begin{flushleft}
		\scriptsize
		\textbf{Estado:}
		\quad
		\textcolor{completado}{\rule{0.45cm}{0.16cm}}
		\; Completado
		\qquad
		\textcolor{proceso}{\rule{0.45cm}{0.16cm}}
		\; En proceso
	\end{flushleft}
\end{center}

\subsection{Plan de ejecución}

El plan de ejecución organiza las acciones de análisis, diseño y evaluación previstas para los tres objetivos específicos de TravelDLS, junto con sus instrumentos de verificación, protagonistas y resultados esperados. La construcción de componentes y las integraciones técnicas se describen en la metodología de desarrollo. Para el primer objetivo se contempla una revisión documental con cadenas de búsqueda en Google Académico, tomando como referencia el mapeo sistemático de \textcite{petersen2008systematic}. La identificación de requisitos se fundamenta en la documentación y los modelos del proyecto; no corresponde a un levantamiento formal con usuarios.

	\begingroup
	\scriptsize
	\setlength{\tabcolsep}{2pt}
	\renewcommand{\arraystretch}{1.15}
	\begin{longtable}{@{}>{\RaggedRight\arraybackslash}p{1.5cm}>{\RaggedRight\arraybackslash}p{2.3cm}>{\RaggedRight\arraybackslash}p{3.75cm}>{\RaggedRight\arraybackslash}p{2.9cm}>{\RaggedRight\arraybackslash}p{2.1cm}>{\RaggedRight\arraybackslash}p{2.85cm}@{}}
	\caption{Plan de ejecución de TravelDLS}\label{tab:plan_ejecucion}\\
	\toprule
	\textbf{Objetivos} & \textbf{Acciones} & \textbf{Descripción de las acciones} & \textbf{Instrumento de verificación} & \textbf{Protagonistas} & \textbf{Resultados esperados} \\
	\midrule
	\endfirsthead
	\caption[]{Plan de ejecución de TravelDLS (continuación)}\\
	\toprule
	\textbf{Objetivos} & \textbf{Acciones} & \textbf{Descripción de las acciones} & \textbf{Instrumento de verificación} & \textbf{Protagonistas} & \textbf{Resultados esperados} \\
	\midrule
	\endhead
	\endfoot
	\bottomrule
	\endlastfoot
	\hypertarget{plan-oe1}{OE1} & Revisión documental y diagnóstico & Seleccionar antecedentes, examinar los procesos de transporte de carga pesada y delimitar necesidades de gestión, seguimiento y planificación de rutas. & Cadenas de búsqueda aplicadas en Google Académico, tomando como referencia el procedimiento de mapeo sistemático (Systematic Mapping) de Petersen et al. (2008). & Equipo investigador. & Problemática fundamentada y necesidades iniciales identificadas. \\ \midrule
	OE1 & Identificación de requisitos del sistema & Analizar los objetivos, el alcance y los modelos de TravelDLS para identificar las funcionalidades y necesidades técnicas de la plataforma. No se realizó levantamiento directo con usuarios. & Documento del proyecto, diagramas UML y entidad-relación. & Equipo investigador. & Se identificaron preliminarmente las funcionalidades y componentes de TravelDLS mediante análisis documental, sin validación directa con usuarios. \\ \midrule
	OE1 & Análisis de métodos de rutas & Definir un escenario hipotético con múltiples destinos, calcular recorridos con Vecino Más Cercano y 2-opt, y estudiar la formulación de Google OR-Tools para una comparación posterior bajo las mismas condiciones. & Matriz de distancias hipotética, cálculos reproducibles y criterios de comparación. & Equipo investigador. & Fundamentos y datos de referencia para evaluar los métodos cuando OR-Tools se integre; sin atribuirle resultados de ejecución. \\ \midrule
	OE2 & Definición de arquitectura & Establecer la organización del frontend con Next.js, el backend con AdonisJS bajo MVC y capa de servicios, la base de datos PostgreSQL y los mecanismos de acceso por empresa. & Documento de arquitectura y esquema de componentes. & Equipo de desarrollo. & Arquitectura y responsabilidades técnicas de los componentes definidas. \\ \midrule
	\hypertarget{plan-oe2}{OE2} & Modelado de datos y clases & Elaborar y revisar los diagramas de clases y entidad-relación, considerando usuarios, roles, empresas, vehículos, conductores, clientes, órdenes y pagos. & Diagramas UML y entidad-relación revisados. & Equipo de desarrollo. & Modelos coherentes con las entidades y relaciones previstas en TravelDLS. \\ \midrule
	OE2 & Modelado de procesos & Revisar los diagramas de casos de uso, estados y secuencia, unificando las reglas sobre cotización, pago, asignación del conductor y seguimiento de las órdenes. & Diagramas de casos de uso, estados y secuencia; registro de revisión. & Equipo de desarrollo y asesoría académica. & Flujos de negocio y responsabilidades consistentes entre los modelos. \\ \midrule
	\hypertarget{plan-oe3}{OE3} & Diseño conceptual del módulo de optimización & Especificar los destinos, costos, restricciones, entradas y salidas del modelo de rutas con Google OR-Tools para orientar su futura integración, sin presentar el módulo como implementado. & Formulación del problema, matriz de distancias y criterios de evaluación. & Equipo investigador y de desarrollo. & Diseño documentado y compromiso de integrar y evaluar el módulo en una etapa posterior. \\ \midrule
	OE3 & Evaluación preliminar del prototipo & Definir escenarios y criterios de prueba, ejecutar los casos disponibles y registrar los resultados observados y los aspectos aún no verificados. & Casos de prueba, registro de ejecución y matriz de resultados. & Equipo de desarrollo y participantes de las pruebas. & Evidencia delimitada de los comportamientos comprobados y de las verificaciones pendientes. \\
	\end{longtable}
	\endgroup

\clearpage
\subsection{Resultados del plan de ejecución}
\label{sec:resultados-plan-ejecucion}

Los resultados del análisis, diseño y evaluación preliminar de TravelDLS se presentan de acuerdo con las acciones establecidas en el plan de ejecución. En relación con el \hyperlink{plan-oe1}{OE1}, se realizó la revisión documental y la identificación preliminar de los requisitos del sistema, lo que permitió fundamentar la problemática e identificar las necesidades iniciales de la plataforma. Para orientar la búsqueda de información, se definieron términos en español e inglés relacionados con el seguimiento y rastreo de vehículos y carga, el transporte de mercancías, la logística, los envíos y la optimización de rutas. Asimismo, se analizaron los objetivos, el alcance y los modelos del proyecto para establecer las funcionalidades y los componentes principales de TravelDLS, sin realizar un levantamiento directo con usuarios.

Las cadenas de búsqueda utilizadas para la revisión documental combinan palabras clave y operadores booleanos, como \texttt{AND} y \texttt{OR}, con el propósito de orientar la consulta de información relacionada con el transporte de carga, el seguimiento de vehículos, la optimización de rutas y las herramientas tecnológicas consideradas para el desarrollo de TravelDLS.

\textbf{Cadenas de búsqueda utilizadas}

\textbf{1. Seguimiento y transporte de carga}

\begin{quote}
	\begin{lstlisting}[style=cadenas]
		("tracking" OR "seguimiento" OR "rastreo" OR "monitoreo") AND ("carga" OR "transporte de carga" OR "freight" OR "freight transportation" OR "logistics")
	\end{lstlisting}
\end{quote}

\begin{quote}
	\begin{lstlisting}[style=cadenas]
		("freight transportation" OR "heavy freight" OR "cargo transportation" OR "logistics") AND ("tracking" OR "real-time tracking" OR "vehicle tracking" OR "cargo tracking") AND ("resource management" OR "fleet management" OR "transport management" OR "logistics management")
	\end{lstlisting}
\end{quote}

Las cadenas de búsqueda de este apartado se utilizaron para localizar investigaciones relacionadas con el seguimiento de cargas, la localización de vehículos y la gestión del transporte. La información obtenida sirvió principalmente para elaborar los antecedentes del proyecto, identificando sistemas e investigaciones que abordan la administración de fletes, el seguimiento en tiempo real y la coordinación de las operaciones logísticas. También contribuyó al marco teórico, especialmente en la definición de conceptos relacionados con la gestión logística, el seguimiento de cargas y la localización de vehículos.

\textbf{2. Optimización de rutas y algoritmos}

\begin{quote}
	\begin{lstlisting}[style=cadenas]
		("traveling salesman problem" OR TSP OR "vehicle routing problem" OR VRP) AND ("nearest neighbor" OR "nearest neighbour" OR "2-opt" OR "OR-Tools") AND ("logistics" OR "freight transportation" OR "delivery routing")
	\end{lstlisting}
\end{quote}


\begin{quote}
	\begin{lstlisting}[style=cadenas]
		("2-opt" OR "two-opt" OR "2 opt") AND ("traveling salesman problem" OR "vehicle routing problem" OR "route optimization") AND ("logistics" OR "transportation" OR "freight transport")
	\end{lstlisting}
\end{quote}
Estas cadenas se utilizaron para buscar investigaciones sobre la planificación y optimización de rutas, con énfasis en el Problema del Agente Viajero (TSP), el algoritmo Vecino Más Cercano y el método 2-opt. Los resultados se emplearon en los antecedentes para identificar trabajos que han aplicado métodos de optimización en el transporte y la distribución de mercancías. Asimismo, sirvieron para desarrollar las bases teóricas del proyecto, explicar los métodos considerados y fundamentar la comparación de alternativas para la planificación de recorridos en TravelDLS.

\textbf{3. Google OR-Tools y matrices de distancias}

\begin{quote}
	\begin{lstlisting}[style=cadenas]
		("Google OR-Tools" OR "OR-Tools" OR "Google Operations Research Tools") AND ("vehicle routing problem" OR "traveling salesman problem" OR "route optimization") AND ("logistics" OR "transportation" OR "delivery")
	\end{lstlisting}
\end{quote}

\begin{quote}
	\begin{lstlisting}[style=cadenas]
		("distance matrix" OR "travel time matrix" OR "road network distances") AND ("route optimization" OR "vehicle routing problem" OR "traveling salesman problem") AND ("logistics" OR "freight transportation" OR "delivery routing")
	\end{lstlisting}
\end{quote}

Las cadenas de este apartado se utilizaron para investigar el funcionamiento y la aplicación de Google OR-Tools en problemas de optimización de rutas, así como el uso de matrices de distancias y tiempos de recorrido. La información recopilada contribuyó principalmente a las bases teóricas y a la fundamentación técnica del módulo de optimización de TravelDLS. Además, permitió identificar elementos necesarios para comparar diferentes recorridos y comprender cómo estas herramientas pueden apoyar la planificación de viajes con múltiples destinos.

\textbf{4. Geolocalización y seguimiento mediante Mapbox}

\begin{quote}
	\begin{lstlisting}[style=cadenas]
		("Mapbox" OR "Mapbox API") AND ("truck tracking" OR "vehicle tracking" OR "fleet tracking" OR "GPS tracking") AND ("logistics" OR "freight transportation" OR "cargo transportation" OR "delivery")
	\end{lstlisting}
\end{quote}

\begin{quote}
	\begin{lstlisting}[style=cadenas]
		("Mapbox" OR "Mapbox API" OR "Mapbox SDK") AND ("real-time tracking" OR "vehicle tracking" OR "GPS tracking" OR geolocation) AND (logistics OR transportation OR "fleet management" OR delivery OR travel)
	\end{lstlisting}
\end{quote}

La consulta de documentación sobre Mapbox, sus API y SDK permitió conocer las herramientas disponibles para trabajar con mapas, geolocalización y seguimiento de vehículos. La información obtenida se consideró en la selección de tecnologías y en el análisis de la viabilidad técnica del proyecto, como parte de la evaluación de las herramientas que podrían apoyar las funcionalidades de ubicación y seguimiento de recorridos.

\textbf{5. Almacenamiento de imágenes mediante Cloudinary}

\begin{quote}
	\begin{lstlisting}[style=cadenas]
		("Cloudinary" OR "Cloudinary API") AND ("web application" OR "mobile application" OR platform OR system) AND (logistics OR transportation OR travel OR delivery OR tracking OR fleet) AND ("image storage" OR "image management" OR "media storage" OR "cloud storage")
	\end{lstlisting}
\end{quote}

Las búsquedas sobre Cloudinary y el almacenamiento de imágenes y recursos multimedia en aplicaciones web permitieron consultar información sobre servicios para gestionar este tipo de archivos. Los resultados se consideraron en la selección de tecnologías y en el análisis de la viabilidad técnica de TravelDLS, con el propósito de evaluar una alternativa para el almacenamiento y la gestión de recursos gráficos asociados al sistema.

\paragraph{Ejercicios preliminares de optimización de rutas}
Para la acción «Análisis de métodos de rutas» del OE1 se prepararon dos ejercicios \emph{hipotéticos y reproducibles}: un camino entre una recogida y una entrega, como ilustración complementaria, y un recorrido con cuatro destinos, relacionado directamente con el escenario del plan. Estos ejercicios permiten documentar los cálculos de Vecino Más Cercano y 2-opt y establecer una matriz común para la comparación futura con Google OR-Tools. Las distancias siguientes pertenecen únicamente a los datos ficticios de cada ejercicio; no son distancias reales ni resultados obtenidos por OR-Tools.

\textbf{Ejercicio complementario: trayecto directo A--B.} Se consideraron cuatro nodos: A (recogida), B (entrega) y dos intersecciones, X e Y. Las conexiones disponibles y sus distancias hipotéticas son A--X (4 km), X--B (6 km), A--Y (3 km), Y--B (10 km) y A--B (12 km). Se supone que pueden recorrerse en los sentidos indicados. La Tabla~\ref{tab:caminos-hipoteticos} compara los tres caminos posibles de esta red simplificada.

\begin{table}[H]
\centering
\caption{Caminos del ejercicio directo A--B}
\label{tab:caminos-hipoteticos}
\begin{tabular}{lcc}
\toprule
Camino & Cálculo & Distancia \\
\midrule
A--X--B & $4+6$ & 10 km \\
A--B & $12$ & 12 km \\
A--Y--B & $3+10$ & 13 km \\
\bottomrule
\end{tabular}
\end{table}

El menor costo es 10 km para A--X--B. Este valor se obtuvo por enumeración de los caminos del ejemplo; queda pendiente comprobar la salida de \texttt{SimpleMinCostFlow} con la misma red. En este caso no hay un orden de entregas que reorganizar, por lo que Vecino Más Cercano y 2-opt no se aplican a la comparación.

\textbf{Ejercicio principal: una recogida y cuatro entregas.} Se fijó A como punto de recogida inicial y B, C, D y E como destinos obligatorios. Para este ejercicio se supone un solo vehículo, distancias simétricas y retorno a A. La Tabla~\ref{tab:matriz-rutas-hipotetica} muestra la matriz ficticia, expresada en kilómetros.

\begin{table}[H]
\centering
\caption{Matriz de distancias hipotéticas para cinco ubicaciones}
\label{tab:matriz-rutas-hipotetica}
\begin{tabular}{crrrrr}
\toprule
Nodo & A & B & C & D & E \\
\midrule
A & 0 & 7 & 8 & 9 & 13 \\
B & 7 & 0 & 1 & 7 & 7 \\
C & 8 & 1 & 0 & 7 & 7 \\
D & 9 & 7 & 7 & 0 & 14 \\
E & 13 & 7 & 7 & 14 & 0 \\
\bottomrule
\end{tabular}
\end{table}

Al aplicar la técnica de Vecino Más Cercano \parencite{contreras2021aplicacion}, desde A se elige B (7 km) y luego C (1 km). Ante el empate de 7 km entre D y E desde C, se adopta el desempate alfabético y se elige D. Con los datos ficticios de la Tabla~\ref{tab:matriz-rutas-hipotetica}, el recorrido resultante A--B--C--D--E--A mide $7+1+7+14+13=42$ km.

Para mejorar ese recorrido se aplica la heurística 2-opt, que intercambia dos aristas e invierte el segmento intermedio \parencite{dacosta2021twoopt,uddin2023twoopt}. En este ejemplo se reemplazan A--B y D--E, cuyo costo conjunto es 21 km, por A--D y B--E, cuyo costo es 16 km. Según la misma matriz ficticia, la inversión produce A--D--C--B--E--A, con una distancia de $9+7+1+7+13=37$ km. La reducción respecto de la ruta inicial es $(42-37)/42\times100=11{,}90$\,\%.

La enumeración de las $4!=24$ secuencias de los destinos, manteniendo fijo A al inicio y al final, confirma que 37 km es el mínimo de esta matriz; la ruta mostrada es una de las que alcanzan ese valor. Esto no demuestra que 2-opt encuentre siempre un mínimo global. La Tabla~\ref{tab:comparacion-rutas-preliminar} separa los cálculos ya comprobados de la evaluación aún pendiente.

\begin{table}[H]
\centering
\caption{Estado de la comparación del ejercicio con múltiples destinos}
\label{tab:comparacion-rutas-preliminar}
\begin{tabular}{@{}p{0.52\textwidth}p{0.18\textwidth}p{0.2\textwidth}@{}}
\toprule
Procedimiento & Distancia & Estado \\
\midrule
Vecino Más Cercano & 42 km & Cálculo manual \\
Vecino Más Cercano y 2-opt & 37 km & Cálculo manual \\
Enumeración de 24 recorridos & 37 km & Mínimo verificado \\
Google OR-Tools (\texttt{RoutingModel}) & Pendiente & Sin implementación \\
\bottomrule
\end{tabular}
\end{table}

Se asume el compromiso de integrar el módulo de optimización con Google OR-Tools en una etapa posterior del desarrollo de TravelDLS, conforme al objetivo específico del proyecto. Para completar la comparación se ejecutará \texttt{RoutingModel} sobre la misma matriz y condiciones, se registrarán el recorrido obtenido y su distancia, y se medirán los tiempos de cálculo bajo un entorno común. La demostración directa podrá comprobarse por separado con \texttt{SimpleMinCostFlow}. Hasta entonces, este resultado documenta una preparación matemática y criterios para la evaluación futura, sin atribuir al prototipo rutas optimizadas ni dar por cumplida la implementación del módulo.

En relación con el \hyperlink{plan-oe2}{OE2}, se definió una arquitectura que separa la interfaz desarrollada con Next.js, los controladores y servicios del backend construido con AdonisJS y la persistencia en PostgreSQL. Se elaboraron los modelos de clases y entidad-relación para representar usuarios, roles, empresas, vehículos, conductores, clientes, órdenes, cargas y pagos, así como sus asociaciones. Estos modelos permiten describir la organización prevista de los datos y las responsabilidades de los componentes. Aún queda pendiente documentar en un esquema de componentes la interacción entre las tres capas y comprobar de manera sistemática la correspondencia de los modelos con las estructuras implementadas y las reglas de acceso por empresa. Por ello, se presenta un diseño elaborado, sin afirmar que su revisión integral se haya completado.

Como evidencia de las acciones de «Modelado de datos y clases» y «Modelado de procesos» previstas para el OE2 en el plan de ejecución, se presentan a continuación los diagramas elaborados. Primero se muestran los modelos de la estructura y los datos del sistema; después, los modelos de sus procesos. Estas figuras documentan el avance del diseño, cuya correspondencia con la implementación todavía requiere una revisión sistemática.

\newcommand{\reservarEspacioFigura}[1]{%
	\par
	\ifdim\dimexpr\pagegoal-\pagetotal\relax<#1\relax
		\newpage
	\fi
}

\reservarEspacioFigura{0.62\textheight}
\begin{samepage}
En la Figura~\ref{fig:diagrama-clases} se presentan las clases principales de TravelDLS, sus atributos y las relaciones entre usuarios, empresas, clientes, conductores y órdenes. También se muestran los servicios encargados de la autenticación, la gestión de órdenes y los pagos simulados.

\begin{figure}[H]
	\centering
	\includegraphics[width=\textwidth,height=0.82\textheight,keepaspectratio]{diagrama-clases.png}
	\caption{Diagrama de clases de TravelDLS: modelado de datos y clases del OE2}
	\label{fig:diagrama-clases}
\end{figure}
\end{samepage}

\medskip
\reservarEspacioFigura{0.78\textheight}
\begin{samepage}
En la Figura~\ref{fig:diagrama-er} se presenta el diagrama entidad-relación de TravelDLS, que muestra las tablas de la base de datos y cómo se relacionan entre sí. Este diagrama permite comprender cómo se almacena la información del sistema y sirve como guía para organizar los datos durante el desarrollo.

\begin{figure}[H]
	\centering
	\includegraphics[width=\textwidth,height=0.82\textheight,keepaspectratio]{diagrama-er.png}
	\caption{Diagrama entidad-relación de TravelDLS: modelado de datos y clases del OE2}
	\label{fig:diagrama-er}
\end{figure}
\end{samepage}

Los diagramas de casos de uso, estados y secuencia corresponden a la acción «Modelado de procesos» del OE2. El primero presenta las funciones por actor; el segundo, los estados de una orden; y el tercero, las interacciones del flujo principal de transporte. En comparación con el resultado esperado de obtener flujos y responsabilidades consistentes entre los modelos, también se considera un cumplimiento parcial: los procesos previstos están representados, pero no se incorpora un registro de revisión comparativa que confirme la coherencia de todas sus reglas.

\medskip
\reservarEspacioFigura{0.84\textheight}
\begin{samepage}
En la Figura~\ref{fig:diagrama-casos-uso} se describen las funciones previstas para el cliente, la empresa transportista y el conductor, como solicitar transporte, gestionar recursos y reportar la ubicación del viaje. El actor «Usuario» agrupa las funciones compartidas de acceso y administración de la cuenta.

\begin{figure}[H]
	\centering
	\includegraphics[width=0.88\textwidth,height=0.82\textheight,keepaspectratio]{diagrama-casos-uso.png}
	\caption{Diagrama de casos de uso de TravelDLS: modelado de procesos del OE2}
	\label{fig:diagrama-casos-uso}
\end{figure}

\smallskip
La Figura~\ref{fig:diagrama-estados} representa los estados de una orden desde su creación hasta la entrega o anulación, incluidos el pago, la asignación del conductor y las transiciones ante fallos.

\begin{figure}[H]
	\centering
	\includegraphics[width=\textwidth,height=0.82\textheight,keepaspectratio]{diagrama-estado-actualizado.png}
	\caption{Diagrama de estados de una orden de transporte: modelado de procesos del OE2}
	\label{fig:diagrama-estados}
\end{figure}
\end{samepage}

\medskip
\reservarEspacioFigura{0.86\textheight}
\begin{samepage}
En la Figura~\ref{fig:diagrama-secuencia} se muestra el intercambio de solicitudes y respuestas entre el cliente, la empresa, el conductor, el frontend, el backend y la base de datos. Las interacciones se organizan en creación del pedido, aprobación del precio, pago simulado, asignación del conductor y seguimiento hasta la entrega.

\begin{figure}[H]
	\centering
	\includegraphics[width=\textwidth,height=0.75\textheight,keepaspectratio]{diagrama-secuencia.png}
	\caption{Diagrama de secuencia del proceso principal de una orden: modelado de procesos del OE2}
	\label{fig:diagrama-secuencia}
\end{figure}
\end{samepage}

\clearpage
En relación con el \hyperlink{plan-oe3}{OE3}, se formuló un diseño conceptual del módulo de optimización para recorridos con múltiples destinos. Se definieron un origen, una matriz de costos y la condición de visitar todos los destinos, así como la necesidad de registrar el orden de visita, la distancia total y el tiempo de cálculo cuando se ejecute la comparación. Este resultado corresponde a la especificación del modelo; Google OR-Tools aún no se integró al prototipo. Se mantiene el compromiso de implementar el módulo y comprobar sus recorridos en una etapa posterior, por lo que el objetivo de implementación no se presenta como cumplido.

Para la acción «Evaluación preliminar del prototipo» se incorporaron pruebas unitarias de reglas de entrega, cifrado, validación de imágenes y operaciones de pago, junto con pruebas funcionales de autenticación, recursos de transporte, órdenes, permisos y respuestas de la API. En una ejecución realizada el 9 de noviembre de 2026 con \texttt{pnpm test} sobre una base de datos aislada para pruebas, se registraron 14 casos unitarios y 33 funcionales aprobados, sin fallos. La Tabla~\ref{tab:pruebas-backend-oe3} resume el resultado observado; el registro completo de la ejecución se conserva como evidencia técnica.

\begin{table}[H]
\centering
\caption{Resultado de la evaluación preliminar del backend para el OE3}
\label{tab:pruebas-backend-oe3}
\begin{tabular}{@{}lrrr@{}}
\toprule
Suite & Aprobadas & Fallidas & Total \\
\midrule
Unitaria & 14 & 0 & 14 \\
Funcional & 33 & 0 & 33 \\
\midrule
Total & 47 & 0 & 47 \\
\bottomrule
\end{tabular}
\end{table}

La aprobación de los casos ejecutados documenta únicamente el comportamiento verificado en esos escenarios; no equivale a una evaluación integral del prototipo. Aún falta consolidar una matriz que relacione cada caso funcional con su resultado, fecha e incidencia, y ampliar las pruebas de flujos completos, navegación e integraciones externas. La evaluación del módulo de optimización también permanece pendiente hasta que se implemente Google OR-Tools.

\subsection{Compromisos para la siguiente etapa}
\label{sec:compromisos-siguiente-etapa}

Los ejercicios preliminares de rutas y la consulta actual de ubicaciones permiten delimitar el trabajo pendiente, pero no acreditan la integración de Google OR-Tools ni de Mapbox. Para completar estas capacidades se asumen los siguientes compromisos:

\begin{itemize}
	\item Se completará el módulo de optimización con Google OR-Tools para recorridos con múltiples destinos. La implementación utilizará la matriz de distancias y las condiciones definidas en los ejercicios preliminares. Después, se comparará el recorrido obtenido con Vecino Más Cercano y 2-opt bajo los mismos datos y se registrarán el orden de visita, la distancia total y el tiempo de cálculo de cada método. Esta comparación permitirá evaluar el módulo con evidencia reproducible, sin anticipar resultados que aún no se han comprobado.
	\item Se integrará Mapbox en la visualización y consulta de ubicaciones del seguimiento de viajes. Actualmente, esas vistas utilizan OpenStreetMap y Leaflet; por ello, será necesario adaptar el mapa para representar con Mapbox la posición reportada por el conductor y asociada a cada orden. Una vez realizada la integración, se comprobará que los participantes autorizados puedan consultar la ubicación correspondiente y que esta se muestre correctamente en la interfaz. Hasta completar esa verificación, Mapbox permanecerá como componente pendiente.
\end{itemize}

% Capítulo 4: Desarrollo
% Propósito: Documentar el análisis de requerimientos, diseño de arquitectura, modelado de base de datos, desarrollo de componentes (backend/frontend) y pruebas del software.

% Capítulo 5: Conclusiones
% Propósito: Sintetizar las conclusiones obtenidas en función de los objetivos alcanzados y plantear recomendaciones o trabajos futuros.

\section{Conclusiones}

El análisis realizado permitió definir los requerimientos iniciales de TravelDLS y establecer una base para evaluar los métodos de optimización de rutas. Mediante un ejercicio hipotético con múltiples destinos, Vecino Más Cercano generó un recorrido de 42 km y 2-opt lo redujo a 37 km, equivalente a una mejora del 11,90 \%. Estos resultados corresponden únicamente a datos ficticios utilizados con fines de comparación; la evaluación con Google OR-Tools bajo las mismas condiciones aún permanece pendiente.

La arquitectura del sistema quedó definida mediante un frontend en Next.js, un backend en AdonisJS y PostgreSQL para la persistencia de datos. Los diagramas de clases, entidad-relación, casos de uso, estados y secuencia permitieron representar las principales entidades, relaciones y flujos de TravelDLS, estableciendo la estructura necesaria para orientar el desarrollo de los módulos de la plataforma.

El prototipo alcanzó avances funcionales en autenticación, gestión de usuarios y recursos de transporte, órdenes, cotizaciones, asignación de conductores, pagos simulados, ubicación y confirmación de entregas. La evaluación preliminar del backend registró 47 pruebas aprobadas y ninguna fallida, correspondientes a 14 pruebas unitarias y 33 funcionales. Google OR-Tools, Mapbox, el pago real y algunas verificaciones de flujos completos permanecen pendientes.

TravelDLS presenta una base funcional suficiente para continuar su desarrollo como plataforma orientada a la gestión, localización y seguimiento del transporte terrestre de carga pesada. El cumplimiento completo del objetivo general dependerá de finalizar las integraciones pendientes y validar de forma integral los principales flujos del sistema.

A partir de las conclusiones obtenidas y de los aspectos que aún requieren desarrollo o validación, se presentan las siguientes recomendaciones con el propósito de orientar las próximas mejoras de TravelDLS y fortalecer la continuidad del proyecto.

% ==========================================
% Referencias Bibliográficas
% ==========================================
\clearpage
\printbibliography[heading=bibintoc, title={Referencias}]

% ==========================================
% Anexos
% ==========================================
\clearpage
\appendix
\phantomsection
\addcontentsline{toc}{section}{Apéndice: Modelo de flujo de caja de TravelDLS}

\section{Modelo de flujo de caja de TravelDLS}
\label{anexo:flujo-caja}

Este anexo presenta el modelo de flujo de caja utilizado para respaldar la viabilidad económica de TravelDLS. La información fue transcrita del archivo digital \href{run:anexos/flujo-caja-traveldls.xlsx}{flujo-caja-traveldls.xlsx}, incluido como anexo del proyecto. Los supuestos base se presentan en la Tabla~\ref{tab:supuestos-economicos}, dentro de la viabilidad económica.

\begin{table}[H]
\centering
\caption{Contribución estimada por envío}
\label{tab:anexo-contribucion-envio}
\begin{tabularx}{\textwidth}{>{\RaggedRight\arraybackslash}X r >{\RaggedRight\arraybackslash}X}
\toprule
Concepto & Valor (C\$) & Explicación \\
\midrule
Precio final que paga el cliente & 3,954.96 & Precio base equivalente a USD~100 más el porcentaje de TravelDLS. \\
Ingreso bruto de TravelDLS & 292.96 & Porcentaje agregado al precio base, correspondiente al 8\,\%. \\
Comisión cobrada por la pasarela & 118.65 & Costo estimado por procesar el pago, equivalente al 3\,\% del total cobrado al cliente. \\
Ingreso neto de TravelDLS por envío & 174.31 & Monto que queda para TravelDLS después de pagar la comisión de la pasarela. \\
\bottomrule
\end{tabularx}
\end{table}

\begin{table}[H]
\centering
\caption{Flujo de caja proyectado}
\label{tab:anexo-flujo-caja-proyectado}
{\small
\setlength{\tabcolsep}{4pt}
\resizebox{\textwidth}{!}{%
\begin{tabular}{lrrrrrr}
\toprule
Concepto & Año 0 & Año 1 & Año 2 & Año 3 & Año 4 & Año 5 \\
\midrule
Envíos por mes & 0 & 60 & 90 & 120 & 150 & 180 \\
Ingresos de TravelDLS & -- & 125,503.20 & 188,254.80 & 251,006.40 & 313,758.00 & 376,509.60 \\
Costos operativos & -- & 68,992.08 & 69,431.52 & 69,431.52 & 69,431.52 & 69,431.52 \\
Inversión inicial & $-$56,834.24 & -- & -- & -- & -- & -- \\
Flujo neto & $-$56,834.24 & 56,511.12 & 118,823.28 & 181,574.88 & 244,326.48 & 307,078.08 \\
Flujo acumulado & $-$56,834.24 & $-$323.12 & 118,500.16 & 300,075.04 & 544,401.52 & 851,479.60 \\
\bottomrule
\end{tabular}%
}
}
\end{table}

\begin{table}[H]
\centering
\caption{Indicadores financieros del modelo}
\label{tab:anexo-indicadores-financieros}
\begin{tabularx}{\textwidth}{l r l >{\RaggedRight\arraybackslash}X}
\toprule
Indicador & Resultado & Unidad & Concepto \\
\midrule
VAN & 547,106.97 & C\$ & Valor actual de todos los flujos futuros menos la inversión inicial, usando una tasa de descuento del 12\,\%. \\
TIR & 162 & \% & Rentabilidad implícita del proyecto según sus flujos de caja. \\
PRI & 1.002 & Años & Tiempo aproximado que tarda en recuperarse la inversión inicial. \\
\bottomrule
\end{tabularx}
\end{table}


\end{document}
